\chapter{A complete company and a real disclosure}
\label{ch:case}

\section{What a balance must remember}

Northstar Components is a synthetic seller of one homogeneous good. Its
disclosures, transactions and contract terms are specified together, so we can
check whether reconstruction preserves a known possible history. All amounts
in this example are currency units, not thousands. The company has no foreign
currency, subsidiaries, leases, returns or other comprehensive income. These
are declared features of the example, not assumptions to impose on an imported
company. Its source package and expected results are supplied in
\code{fixtures/northstar.json}.

At the beginning of 2024, Northstar holds cash of 60, receivables of 20,
inventory of 30 and plant costing 100 with accumulated depreciation of 20.
It owes suppliers 20 and lenders 100. Share capital is 50 and retained earnings
are 20. Thus assets of 190 equal liabilities of 120 plus equity of 70.
Inventory consists of six units costing 5 each. The existing plant depreciates
by 10 a year, with no residual value. The loan remains outstanding throughout
2024; both possible contract types charge 5\% during that year.

The six months ending June and December provide two successive observation
periods. Here the chosen accrual convention assigns one half-year to each.
Cash-flow dates within a half-year are not recovered from a period aggregate;
the synthetic history supplies a feasible ordering in which inventory is
purchased before delivery and customer receipts precede the payments they fund.
More detailed liquidity questions require more detailed dates.

\begin{center}
\begin{tabular}{lrr}
\toprule Movement & First half & Second half\\\midrule
Delivered sales on credit &110&150\\
Customer collections &110&135\\
Inventory purchases on credit &70&85\\
Supplier payments &65&80\\
Cost of goods delivered &75&90\\
Payroll paid &20&22\\
Plant purchased at period end &10&12\\
Depreciation &5&6\\
Interest accrued and paid &2.5&2.5\\
Tax expense &1.5&5.9\\
Tax paid &3.5&3.9\\
Dividends paid &4&5\\\bottomrule
\end{tabular}
\end{center}

Each row has an economic meaning before it has a posting. Delivery of the
good transfers control and creates an unconditional receivable; invoicing is
assumed simultaneous in this example. Collection settles that receivable.
Purchases increase inventory and trade payables; paying the supplier settles
the payable. Delivery separately removes inventory at cost. Plant purchases
are capital expenditure, while depreciation consumes its service potential.
The addition of 10 enters service after June and has a five-year life, producing
one unit of second-half depreciation. The addition of 12 enters service after
December and has a twelve-year life.

Tax is deliberately simple: 20\% of positive profit before tax, with all these
expenses deductible and no deferred tax. The first-half payment exceeds its
tax expense by 2, leaving prepaid tax; that prepayment is used in the second
half. A genuine corporate tax package would need jurisdiction, tax basis,
losses, timing differences and payment schedules. The example's tax rule
governs its own numbers only.

\section{Recovering movements without inventing invoices}

Suppose the annual report discloses sales of 260, cost of sales of 165 and the
opening and closing working-capital balances, while the collection and purchase
totals are hidden. Under Northstar's explicit no-write-off, no-acquisition and
no-currency-effect scope, the three bridges are
\begin{align}
35&=20+260-C,\label{eq:ns-ar}\\
20&=30+P-165,\\
30&=20+P-S.
\end{align}
Here $C$ denotes collections, $P$ purchases and $S$ supplier payments, all in
currency units over the same year. Substitution gives $C=245$, $P=155$ and
$S=145$. These are identifiable aggregate movements within this model. They do
not identify which invoice was paid on which day. If acquisitions or write-offs
are admitted, their movement variables enter the same bridges and these totals
may cease to be unique.

The loan remains different. Write $F$ and $V$ for fixed and floating principal.
The source imposes $F+V=100$, $F,V\geq0$ and
$0.05F+0.05V=5$. The last equality adds no information about the split.
Both endpoints and every intermediate allocation survive. Northstar therefore
has identified working-capital totals and an unidentified financing exposure
in the same joint historical family. A solver must not turn its convenient
witness $F=0$ into the company's recovered contract.

The model's historical slots are one delivery, collection, purchase, supplier
settlement, payroll, plant addition, depreciation, interest, tax payment and
dividend per half-year. Each slot represents a classified aggregate event.
The interim filing supplies the information for allocating the annual movements
between the two half-years; the annual bridges alone identify only their totals.
The template for collection contributes $+C$ to cash and $-C$ to receivables;
its state effect reduces the unpaid receivable pool by $C$. The purchase
template adds inventory quantity $P/5$, inventory cost $P$ and a payable $P$.
A cash-balance identity by itself cannot substitute for either template.

\section{The complete historical statements}

The corrected statements follow from those movements. Temporary income and
expense accounts close to retained earnings at each half-year boundary;
annual income is the sum of the two periods' classified movements, excluding
closing entries.

\begin{center}
\begin{tabular}{lrrr}
\toprule Balance & Opening & June & December\\\midrule
Cash &60&65&74.6\\
Receivables &20&20&35\\
Inventory &30&25&20\\
Plant, net &80&85&91\\
Prepaid tax &0&2&0\\
Total assets &190&197&220.6\\\midrule
Trade payables &20&25&30\\
Debt &100&100&100\\
Share capital &50&50&50\\
Retained earnings &20&22&40.6\\
Liabilities and equity &190&197&220.6\\\bottomrule
\end{tabular}
\end{center}

Annual revenue of 260 less cost of sales 165, payroll 42, depreciation 11 and
interest 5 gives profit before tax of 37. Tax expense is 7.4 and profit is 29.6.
Other comprehensive income is zero under the declared scope, so total
comprehensive income is also 29.6. Equity moves from 70 to
$70+29.6-9=90.6$ after dividends.

The cash-flow statement supplies an independent reconciliation. Collections
245 less supplier payments 145, payroll 42, interest 5 and tax payments 7.4
produce operating cash of 45.6. Investing cash is $-22$ and financing cash
is $-9$, giving a cash increase of 14.6. The indirect route gives the same
operating cash: profit 29.6 plus depreciation 11, minus the receivables
increase 15, plus the inventory decrease 10 and the payables increase 10.
The prepaid-tax balance is zero at both annual endpoints. The two cash routes
use different information and therefore provide a useful check on classification.

\section{A correction changes what was knowable}

The original interim filing reported revenue of 120 and receivables of 30.
Ten units of revenue had been recorded before delivery. Its tax expense was
3.5, profit 14 and equity 80. The annual filing subsequently corrected the
interim revenue and receivable by 10 and recognised the excess tax payment as
prepaid tax of 2. Corrected interim profit is 6 and equity is 72, as above.
The original statement remains an accurate record of what was published, even
though its revenue recognition was wrong under the stated delivery rule.

For the temporal fixture, the original report becomes available on 1 August
2024 and is ingested on 3 August. The correction and annual report become
available on 2 March 2025 and are ingested on 5 March. An investor-information
query at 1 September 2024 00:00 UTC returns 120. For this synthetic example,
take UTC as the reporting timezone and use an instant cutoff of
3 March 2025 at 00:00 UTC, after the correction's entire publication-date
interval has ended. That query returns 110 when evaluated using the subsequently acquired archive. A query
that additionally asks what the database contained on 4 March still returns
120. Publication and ingestion answer different questions.

The annual total 260 and corrected interim total 110 share revision group
\code{v2}, so second-half revenue is $260-110=150$. Subtracting the original
120 would yield 140, but that result mixes incompatible revisions. The query
must reject the subtraction unless a reviewed compatibility relationship
supports it. A later date alone does not establish such a relationship.

\section{One economic scenario, several complete futures}

For 2025, assume demand for 30 units at price 10, unit purchase cost 5 and a
target closing inventory of eight units. Opening inventory is four units, so
the purchase decision is $30+8-4=34$ units, costing 170. Deliveries cost 150,
leaving inventory of 40. Collect 90\% of current sales and 10 of old
receivables: $270+10=280$, leaving receivables of 55. Pay old suppliers 30
and 130 of the new purchases, leaving payables of 40. Payroll is 50, capital
expenditure of 20 occurs at year-end and dividends are 8.

These decisions are inputs to a conditional forecast. They are not estimates
of demand or optimal management responses. No new borrowing is available.
Under the existing asset schedules, depreciation is $10+10/5+12/12=13$.
The new year-end asset begins depreciation in the next period. Let
$\alpha=F/100$. With the floating rate resetting to 7\%, interest is
$I=7-2\alpha$. Profit before tax is
$300-150-50-13-I=80+2\alpha$, giving tax $16+0.4\alpha$ and profit
$64+1.6\alpha$.

\begin{center}
\begin{tabular}{lrr}
\toprule 2025 statement amount & All fixed & All floating\\\midrule
Revenue &300&300\\
Interest expense &5&7\\
Tax expense &16.4&16\\
Profit and comprehensive income &65.6&64\\
Operating cash flow &48.6&47\\
Investing cash flow &$-20$&$-20$\\
Financing cash flow &$-8$&$-8$\\\midrule
Closing cash &95.2&93.6\\
Receivables / inventory / net plant &55 / 40 / 98&55 / 40 / 98\\
Assets &288.2&286.6\\
Payables / debt &40 / 100&40 / 100\\
Equity &148.2&146.6\\\bottomrule
\end{tabular}
\end{center}

The joint symbolic result is cash $93.6+1.6\alpha$, profit
$64+1.6\alpha$ and equity $146.6+1.6\alpha$, with the same
$\alpha\in[0,1]$ throughout. Selecting the smallest cash and largest equity
independently would describe no feasible company. A requested distribution
requires a measure over that same historical family; the accounting equations
do not supply one.

The later method schedules use the floating share $x$ rather than the fixed
share $\alpha$ used here: $x=1-\alpha$. Both parameterisations describe the
same joint family. Its 2025 opening cash is 74.6 and equity is 90.6; the
forecast spends 20 on period-end capital expenditure, distributes 8 and settles
its stipulated tax in the same year. These teaching assumptions are retained
when the later lease example introduces a different tax-payment schedule.

\section{Keeping the same company through every calculation}
\label{sec:northstar-portable}

The first inventory purchase shows what a reproducible company calculation
must retain. Before the purchase, Northstar has six units costing 30 and owes
suppliers 20. It buys fourteen more units for 70 on credit. The journal debits
inventory and credits trade payables by 70; the resulting inventory is twenty
units costing 100, payables are 90, and cash remains 60. A balanced journal
accompanied by twenty-one units of inventory is wrong under the stipulated
unit cost of 5. The quantity change and the monetary posting must therefore be
checked against the same purchase and opening state.

Delivery follows purchase. Its four journal lines recognise the receivable and
revenue together with cost of sales and the reduction in inventory. For the
first half-year, delivery of fifteen units leaves five units costing 25.
Separating the revenue and inventory effects into independently committable
events could expose a state in which the customer has the goods but the seller
still reports them. Treating delivery as one economic event prevents that
intermediate inconsistency. Collections then precede the payments they fund.
This order establishes a feasible aggregate history; it does not recover the
dates of individual invoices.

The portable company example retains twenty-two successive states and journal
steps: ten economic events and a closing entry in each half-year. Each step
identifies its prior state, economic inputs, accounting template and resulting
state. The corrected June statements and December statements refer to this same
chain. The December closing state, with cash 74.6 and equity 90.6, is also the
explicit opening of the 2025 projection. Consequently, changing a historical
purchase or tax payment cannot leave the future opening silently unchanged.
These are proposed histories; publishing them as candidates does not post them
to an accounting book.

The rate comparison checks a different connection. Recall that $x$ is the
floating share of the unchanged principal of 100. For an imposed floating rate
$r$, the fixed part earns interest at 5\%, so annual interest is
\begin{equation}
 I(x;r)=100(1-x)0.05+100xr=5+(100r-5)x.
 \label{eq:northstar-portable-rate}
\end{equation}
At 7\%, increasing $x$ from zero to one adds 2 to interest; at 9\%, it adds 4.
Under the stipulated 20\% tax rule, profit and cash therefore fall by 1.6 or
3.2 respectively. The second rate supplies a useful control: two authoring
frontends might agree because they share an erroneous constant. Independent
arithmetic must recover the changed slope from the changed rate. The reference
exchanges retain the exact expression for every $x\in[0,1]$, including its
dependence across cash, profit, assets and equity. They attach no probabilities
to the endpoints or to intermediate debt allocations.

The opening also determines what kind of forecast this is. The corrected
December and interim information becomes public in March 2025, after the
January origin of the scenario. Using it is an archival conditional exercise:
given that corrected opening, what follows from these imposed decisions and
rates? It cannot establish what an investor could have forecast in January.
An investor-information request at that earlier cutoff rejects the later
disclosure dependency. Separately, a management policy cannot read a future
realised collection merely because the simulation scheduler already stores it.
An imposed scenario assumption and a later observed outcome have different
information meanings.

Partial information must survive the same chain. If the fraction of current
sales collected is missing, delivery can still determine revenue and cost of
sales. Cash, receivables and the ability to fund the proposed settlements remain
unresolved. The partial report identifies its affected lines and retains the
full debt family. If the tax method is unavailable, profit after tax, retained
earnings and the final closing entry are unresolved as well. A zero in a
temporary account would incorrectly suggest that closing had been completed.
The portable values therefore carry an explicit reason and its outstanding
obligation wherever the calculation is unavailable.

There is a similar distinction in reading a note. The separate authored phrase
``substantially all customers pay promptly'' does not imply a 90\% threshold.
An analyst may impose an all-payers assumption for a conditional calculation,
but its decision must preserve the unquantified source reading. Likewise,
interpreting a displayed 100 as nearest-unit rounding requires an explicit
rounding assumption. With ties rounded upward, that assumption permits values
in $[99.5,100.5)$; absent the assumption, the display alone supplies no such
interval. Neither decision grants permission to post a journal.

Finally, a correct purchase must survive an uncertain acknowledgement. The
request identifies Northstar, its book and branch, and the particular purchase;
it also fixes the opening state, expected revision and all input versions.
Suppose the client sends the request but receives no response. If the store
committed nothing, an exact retry may commit the purchase. If the store already
committed it, the same retry returns the original receipt without another
purchase. The client cannot distinguish these possibilities from a timeout.

Receipt lookup therefore precedes comparison with today's head, after the
original request has passed its input and authorization checks. A later,
separately identified purchase may have advanced the book, but it does not
erase the first receipt. Replaying the first request still returns that receipt.
Changing its request identifier or budget changes its full fingerprint and
conflicts with the already recorded business event. A new event presented
against an obsolete opening is instead rejected as stale. Every rejection
leaves the committed states, journals, quantity changes and receipt index
unchanged.

Review decisions, saved search progress and imports have their own owners and
receipts. Correcting the reported revenue invalidates dependent current uses,
including a selection that depended on the old set of eligible readings; it
does not rewrite committed purchases. Saving one searched region of
$[0,1]$ preserves the other region and the next position to visit. Resuming
with a larger budget updates that saved work, while a changed model requires
a different calculation. These distinctions make the complete company example
reconstructible from its inputs and record references. The operation-specific
fields and exact documentary exchanges are described in
Section~\ref{sec:corporate-exchange-reference}; they still require implementation
and actual fault testing before any claim of database durability.


\section{Apple: source labels do not identify exposure}
\label{sec:apple}

The real reference package is Apple's 2024 Form 10-K and first-quarter 2025
Form 10-Q, for fiscal dates 28 September and 28 December 2024 respectively
\citep{apple2024,apple2025q1}. The annual balance sheet also supplies
30 September 2023. The annual filing is archived as the company's PDF from
AnnualReports; the interim filing is archived from SEC EDGAR with its inline
XBRL. Apple's Note 1 describes a 52/53-week year ending on the last Saturday
of September. Calendar-year assumptions would therefore misstate the periods.

The package inventory preserves every annual PDF page and every interim
top-level section, all thirteen annual notes and all ten interim notes.
It records internal cross-references and externally incorporated material,
including the proxy statement, as distinct obligations. Statement and debt
facts below are structured; policies, other notes and exhibits retain their
source locations and specific modelling obligations. This is complete
inventory coverage of the acquired files, with partial semantic formalisation.
It does not claim to reconstruct every Apple contract or to have acquired
every incorporated external document.

Apple's annual balance sheet reports assets 364,980, liabilities 308,030 and
equity 56,950, in millions of US dollars. Note 9 reports term-debt principal
97,341, unamortised discount/premium and issuance costs of $-321$, and hedge
accounting adjustments of $-358$. Thus its carrying amount is
$97{,}341-321-358=96{,}662$, split into current 10,912 and non-current
85,750. Fair value is separately disclosed as 88.4 billion. None of these
three measurements may replace another merely because a table calls it debt.

Note 9 identifies the notes as fixed-rate, then states that interest-rate
swaps effectively convert a portion to floating rates. Note 4 reports
interest-rate contract notional of 14,575 at the annual date. A notional is
neither a liability carrying amount nor sufficient information for an exact
interest sensitivity: reset timing, remaining periods, curve conventions and
hedge coverage still matter. The cross-reference connects the derivative to
the borrowing without erasing either contract.

At the interim date, total term-debt carrying amount is
$10{,}848+83{,}956=94{,}804$. Cash is 30,299, total assets 344,085,
liabilities 277,327 and equity 66,758. The interim debt note's rounded
``94.8 billion'' and the statement's 94,804 million can coexist; their
precision differs. Treating them as conflicting exact numbers would create
an artificial infeasibility.

These balances leave two practical financing questions open: how much cash can
fund a payment, and when the payment falls due. The annual cash and cash
equivalents balance of 29,943 includes 2.6 billion held in escrow and restricted
from general use for the State Aid tax obligation. Note 4 states the inclusion
explicitly, and Note 7 connects the escrow to that obligation
\citep[Notes 4 and 7, printed pp.~37 and 39]{apple2024}. Subtracting the reported
restricted component gives $29{,}943-2{,}600=27{,}343$ million as an indicative
remainder at the displayed amounts. The restricted component is rounded to
tenths of a billion, so this subtraction does not establish an exact freely
available balance. The earlier case annotation that described the reported
cash balance as excluding restricted cash was wrong relative to this footnote.
Correcting that interpretation changes funding calculations without changing
the issuer's reported number. It also requires revisiting any liquidity result
that relied on the exclusion.

The debt movement supplies a second reason to preserve a calculation's missing
terms. The first-quarter cash flow statement reports repayments of term debt of
1,009 million, while carrying amount falls by
$96{,}662-94{,}804=1{,}858$ million. At the displayed amounts, the reconciliation is
\[
       94{,}804=96{,}662-1{,}009-849.
\]
The 849 million is the residual after the reported cash repayment; it does not
identify an exchange loss, an income-statement charge or a particular journal.
Explaining it requires the missing movement information, including applicable
currency, discount/premium and hedge-accounting effects. The quarterly cash
flow statement still reconciles independently: operating inflow 29,935 plus
investing inflow 9,792 less financing outflow 39,371 gives the reported cash
increase of 356, from 29,943 to 30,299 million
\citep[statement of cash flows and Note 6]{apple2025q1}. A balanced cash bridge
therefore does not identify the debt history.

The annual filing's securities exhibit adds terms that a debt-total query
cannot supply. Exhibit 4.1 describes two euro-denominated series bearing the
year 2025 in their names. Its reported outstanding principal is dated
18 October 2024, after the annual balance-sheet date. The description concerns
Apple Inc. as issuer, rather than every subsidiary in the consolidated group
\citep[Exhibit 4.1, pp.~3--4]{apple2024}.

\begin{center}
\begin{tabular}{lrrl}
\toprule
Series & Principal (EUR m) & Fiscal year & Scheduled maturity\\\midrule
0.875\% 2025 Notes &1,250&2025&24 May 2025\\
0.000\% 2025 Notes &1,000&2026&15 November 2025\\\bottomrule
\end{tabular}
\end{center}

Apple's disclosed 52-week fiscal 2025 ends on 27 September 2025 under its
last-Saturday convention. The November maturity therefore belongs to fiscal
2026 even though the bond's name says 2025. The annual Note 9 principal-payment
schedule uses fiscal years, whereas the instrument's stated maturity is a
calendar date. Joining those fields on their year number would put a billion
euros of scheduled principal into the wrong fiscal bucket. Neither the later
holding date nor the euro principal may be silently substituted for a
September US-dollar consolidated balance.

Suppose the disclosed faces remain outstanding until scheduled maturity, with
no early redemption, default or additional withholding-related payment, and
euro remains the payment currency. The annual coupon on the 0.875\% series is
then the reported face multiplied by the stated rate:
\[
       1{,}250\times0.00875=10.9375\quad\text{million EUR}.
\]
Its final scheduled principal plus annual coupon is consequently
1,260.9375 million EUR; the zero-coupon series has 1,000 million EUR of
scheduled principal and zero stated coupon. Both scheduled maturity dates
fall on a Saturday. Exact settlement requires the applicable payment-adjustment
provision and business calendar. Exhibit 4.1 is expressly a summary qualified
by the 2013 and 2018 indentures; its payment-currency exceptions and redemption
rights also matter. The retained 2021 indenture cannot replace those earlier
contracts. The calculation establishes conditional scheduled amounts, while
their exact settlement dates and subsequent holdings remain unresolved.

Those missing facts have different consequences for the financial statements.
Principal and coupon payments reduce cash, but settling principal differs from
recognising interest expense. The latter also depends on carrying amount,
effective-interest accrual and the treatment of discounts and hedges; Note 9
expressly includes applicable amortisation and hedging adjustments in effective
rates. Exchange rates, hedge designations and tax treatment can affect profit,
other comprehensive income, retained earnings and closing debt differently.
The cash flow statement additionally needs the payment classification and any
associated hedge settlements. A usable answer thus carries the conditional
euro schedule together with these unresolved dependencies. A dollar company
forecast must supply consistent holdings, exchange, hedge and tax assumptions
across all affected statements.

To demonstrate a source-to-forecast calculation, freeze the interim balance
sheet and apply a \emph{marginal} financing stress for a following 13-week
period: a 200-basis-point rate increase, no principal change, no offsetting
hedge beyond the explicitly selected exposure, immediate cash settlement and
no tax effect. An assumed effective floating exposure $E=14{,}575$ million
produces $0.02\times(13/52)E=72.875$ million additional interest. This exposure
uses the annual notional as a conditional input; it is not asserted to be the
interim net sensitivity. Under that condition, debit interest expense and
credit cash by 72.875. The resulting marginal statements have cash 30,226.125,
assets 344,012.125, liabilities 277,327, equity 66,685.125, profit and comprehensive
income of $-72.875$, and operating cash flow of $-72.875$ relative to the frozen
baseline. The balance-sheet difference equals the equity difference exactly.

If exposure is instead only restricted to $0\leq E\leq14{,}575$, the
conditional added interest lies in $[0,72.875]$ before monetary rounding. That
restriction is a stress assumption and is not claimed to bound every real
Apple exposure. This real case verifies evidence, context, mappings and a
coherent marginal propagation. Northstar supplies the complete-company
forecast; a complete Apple forecast remains dependent on its unmodelled
operations and undisclosed terms.

\section{Apple: five statements constrain one year}
\label{sec:apple-statements}

The debt example leaves a larger question: can the reported profit, changes
in equity and cash flows describe the same financial year? Apple's five
consolidated statements supply a useful starting point. Their fiscal 2024
amounts refer to the year ending 28 September 2024; the opening balance sheet
is dated 30 September 2023. All monetary amounts below are millions of US
dollars. The operations statement's per-share amounts and weighted share
counts keep their separate units \citep[printed pp.~29--33]{apple2024}.

Products revenue of 294,866 and services revenue of 96,169 sum to 391,035.
Their costs, 185,233 and 25,119, leave gross margin of 180,683. Research and
development of 31,370 and selling, general and administrative expense of
26,097 reduce that amount to operating income of 123,216. Other income of
269 and an income-tax provision of 29,749 give net income of 93,736
\citep[printed p.~29]{apple2024}. This same net income enters both the equity
statement and the operating-cash reconciliation. It is one period result
appearing in three statements, so it must not be added three times.

Other comprehensive income supplies the next connection. Foreign-currency
translation contributes 395 after tax. Derivative fair-value changes of
$-832$ and reclassification adjustments of $-1{,}337$ contribute $-2{,}169$;
marketable-debt fair-value changes of 5,850 and reclassifications of 204
contribute 6,054. Thus OCI is $395-2{,}169+6{,}054=4{,}280$, and total
comprehensive income is $93{,}736+4{,}280=98{,}016$
\citep[printed p.~30]{apple2024}. Adding the period OCI to opening accumulated
other comprehensive loss of $-11{,}452$ gives the closing $-7{,}172$.
That accumulated balance, rather than the period flow, enters the balance
sheet.

Despite positive comprehensive income, total equity falls. The equity
statement explains the difference: common-stock issuance adds 1,423 and
share-based compensation adds 12,034; withholding reduces the capital
component by 3,993 and retained earnings by 1,612; declarations reduce
retained earnings by 15,218 and repurchases by 95,846
\citep[printed p.~32]{apple2024}. These movements, apart from comprehensive
income, total $-103{,}212$. The complete equity reconciliation is therefore
\[
62{,}146+98{,}016+1{,}423+12{,}034-3{,}993-1{,}612
 -15{,}218-95{,}846=56{,}950.
\]
The components reconcile separately: capital rises from 73,812 to 83,276,
the accumulated deficit moves from $-214$ to $-19{,}154$, and accumulated
other comprehensive loss moves from $-11{,}452$ to $-7{,}172$. These are
the three closing equity components reported on the balance sheet. Their
agreement prevents a model from explaining the correct total with the wrong
allocation among capital, retained earnings and accumulated OCI.

The same change appears in net assets. Table~\ref{tab:apple-statement-join}
compares the reported endpoints; its changes are derived subtractions
\citep[printed p.~31]{apple2024}. Assets rise by 12,397 while liabilities
rise by 17,593, leaving a 5,196 decrease in equity, exactly the decrease
from the equity reconciliation.

\begin{table}[htbp]
\centering
\begin{tabular}{lrrr}
\toprule &30 September 2023&28 September 2024&Change\\\midrule
Assets&352,583&364,980&12,397\\
Liabilities&290,437&308,030&17,593\\
Equity&62,146&56,950&$-5{,}196$\\\bottomrule
\end{tabular}
\caption{The balance-sheet change agrees with the complete equity movement.
The agreement constrains a history without identifying its transactions.}
\label{tab:apple-statement-join}
\end{table}

Cash follows a different reconciliation. Starting from the same net income,
the cash-flow statement adds depreciation and amortisation of 11,445,
share-based compensation expense of 11,688 and other adjustments of
$-2{,}266$. These three adjustments sum to 20,867. Its six operating-asset
and liability adjustments are net receivables $-3{,}788$, vendor non-trade
receivables $-1{,}356$, inventory $-1{,}046$, other assets $-11{,}731$,
payables $+6{,}020$ and other liabilities $+15{,}552$. They sum to 3,651,
giving operating cash of $93{,}736+20{,}867+3{,}651=118{,}254$
\citep[printed p.~33]{apple2024}. The full indirect bridge is thus present;
none of its broad adjustment rows has been replaced by a presumed change in
a similarly named balance-sheet account.

Investing cash is also positive: security purchases of $-48{,}656$,
maturities of 51,211, sales of 11,135, plant-acquisition payments of
$-9{,}447$ and other investing cash of $-1{,}308$ sum to 2,935. Financing
uses 121,983: employee-award tax payments of 5,441, dividend payments of
15,234, share repurchases of 94,949, debt repayments of 9,958 and other
outflows of 361 outweigh net commercial-paper proceeds of 3,960. The source
shows a dash for net term-debt issuance; the displayed total treats that row
as making no contribution, while the source reading retains the dash.
The three activity totals give
\[
118{,}254+2{,}935-121{,}983=-794,
\qquad 30{,}737-794=29{,}943.
\]
The opening and closing amounts here are cash, cash equivalents and restricted
cash and cash equivalents, as defined by the cash-flow statement. Opening
balance-sheet cash is only 29,965, so the cash-flow population exceeds it by
772; both closing amounts are 29,943. Consequently, the balance-sheet line
falls by 22 while the cash-flow total falls by 794. The earlier 772 remains
unallocated in the retained evidence. It cannot be relabelled as the closing
2.6 billion escrow, which Note 4 explicitly includes in reported cash and
equivalents \citep[printed p.~33 and Note 4, printed p.~37]{apple2024}.

These joins expose the missing transaction links more precisely. Annual
receivables rise by 3,902, leaving 114 after adding their $-3{,}788$ indirect
adjustment. Inventory rises by 955, leaving $-91$ after its $-1{,}046$
adjustment. Payables rise by 6,349, exceeding their 6,020 adjustment by 329.
Even the exact opposite vendor-receivable movements identify neither gross
purchases nor collections. The notes examined below must supply the relevant
recognition, valuation, currency and settlement information; the residuals
themselves supply no such explanation.

Nor can the equity rows simply replace financing cash flows. The reported
repurchase reduction of 95,846 exceeds cash repurchases of 94,949 by 897;
withholding reductions total 5,605 against cash tax payments of 5,441; and
equity compensation of 12,034 differs from the 11,688 expense adjustment.
Cash dividends of 15,234 exceed declarations of 15,218 by 16. Those
differences remain unassigned to particular settlements, fees, taxes or
recognition periods \citep[printed pp.~32--33]{apple2024}. The later note
examples retain their own evidence and conditional histories.

A reconstruction must now satisfy the statement relationships together.
Correct total cash cannot excuse wrong equity components, and correct net
assets cannot excuse omitted cash-flow adjustments. The five statements
constrain a common history, but their agreement does not reveal all gross
events, dates, counterparties or accounting treatments. Those unresolved
links determine which historical alternatives can be supported and what
additional facts a conditional future path still needs.

\subsection{The following quarter must start from the same balances}

The next filing makes the time connection testable. Its balance-sheet
comparison is 28 September 2024, and all 27 monetary rows reproduce the
annual closing column. Its income, comprehensive-income, equity and cash-flow
comparisons instead cover the three months ending 30 December 2023. Those
earlier-quarter flows cannot become the opening amounts for the quarter ending
28 December 2024. The interim statements are explicitly unaudited
\citep[printed pp.~1--5]{apple2025q1}.

Quarterly products revenue of 97,960 and services revenue of 26,340 total
124,300. Cost of sales of 66,025 and operating expenses of 15,443 leave
operating income of 42,832; other expense of 248 and tax expense of 6,254
give net income of 36,330. Translation OCI of $-625$, derivative OCI of
2,435 and marketable-debt OCI of $-1{,}427$ add 383, giving comprehensive
income of 36,713 \citep[printed pp.~1--2]{apple2025q1}. The resulting
equity movement begins at the same 56,950 that closed the annual record.
Compensation adds 3,383, withholding removes 1,891 from capital and 1,102
from retained earnings, declarations remove 3,819 and repurchases remove
23,476. Together these other movements are $-26{,}905$, so
$56{,}950+36{,}713-26{,}905=66{,}758$. Capital, accumulated deficit and
accumulated OCI also join separately at both endpoints
\citep[printed pp.~3--4]{apple2025q1}.

Assets decline by 20,895 to 344,085 and liabilities decline by 30,703 to
277,327. Their difference gives the same 9,808 equity increase. This check
connects the period result to the balance sheet without confusing declining
assets with a loss. The retained earnings deficit improves by 7,933, capital
rises by 1,492 and accumulated OCI rises by 383; these three changes also
sum to 9,808.

The cash reconciliation supplies a different account of the quarter.
Depreciation and amortisation of 3,080, compensation expense of 3,286 and
other adjustments of $-2{,}009$ add 4,357 to net income. The six operating
asset/liability adjustments are $+3{,}597$, $+3{,}166$, $+215$, $+939$,
$-6{,}671$ and $-11{,}998$, in the same reported categories used for the
annual reconciliation. They total $-10{,}752$, so operating cash is
$36{,}330+4{,}357-10{,}752=29{,}935$. Investing cash adds 9,792 and
financing uses 39,371, leaving a 356 increase from 29,943 to 30,299
\citep[printed p.~5]{apple2025q1}. Both cash-flow endpoints now match the
corresponding balance-sheet amounts. That numeric match still does not make
every dollar unrestricted.

Adjacent flows may now be combined for a deliberately longer question.
Over fiscal 2024 and the following quarter together, comprehensive income is
134,729 and the other equity movements total $-130{,}117$. Over that same
interval, operating, investing and financing cash total 148,189, 12,727 and
$-161{,}354$. The two cumulative reconciliations are
\[
62{,}146+134{,}729-130{,}117=66{,}758,
\qquad
30{,}737+148{,}189+12{,}727-161{,}354=30{,}299.
\]
These are derived amounts for the interval from 1 October 2023 through
28 December 2024, not an issuer annual statement or a trailing-twelve-month
result. Endpoint balances are carried forward, not added. EPS and weighted
share averages also cannot be summed as if they were dollar flows.

The quarter preserves several distinctions that a common schema might
otherwise erase. The annual equity table contains a common-stock issuance
row and the annual cash-flow table contains a term-debt issuance row; the
interim tables omit both. The interim reported totals reconcile over their
own disclosed rows. A comparison record may describe those omissions, but
must not manufacture two issuer-reported zeros. The prior-year quarter's
term-debt repayment dash is likewise retained in its own dated column
\citep[printed pp.~4--5]{apple2025q1}.

Transaction uncertainty remains across the joined interval. In the current
quarter, the receivable change plus its indirect adjustment leaves $-174$;
inventory leaves $-160$. The payable change is $-7{,}050$, differing from
the $-6{,}671$ adjustment by $-379$. Cash dividends exceed declarations by
37, and cash repurchases exceed the equity reduction by 130. None identifies
the missing settlement, valuation or timing components. The subsequent note
inquiries must retain these relationships together when testing historical
alternatives. Combining individually compatible note examples does not by
itself establish that they share a compatible cash, income and equity history.


\section{Apple: revenue recognition and cash collection}
\label{sec:apple-revenue}

An iPhone sale can create revenue today and an obligation to provide services
later. Apple's annual revenue note separates delivered hardware and bundled
software from the right to receive product-related services such as iCloud,
Siri and Maps. It allocates the price and related discounts by relative
stand-alone selling prices. For the bundled services those prices are estimates,
because observable separate prices are unavailable. The hardware portion is
recognised when control transfers, generally on shipment; the service portion
is deferred and recognised straight-line over the estimated period of provision
\citep[Note 2, printed p.~35]{apple2024}. A historical model consequently needs
the allocation and service cohort as well as the sale date and amount. Neither
the product label nor the customer payment supplies those missing terms.

The note describes collection shortly after control transfer or the start of
service delivery. It does not specify one collection lag for every contract.
Products revenue also reflects expected returns, price protection and customer
incentives. Revenue excludes taxes collected for remittance to governments
\citep[Notes 1--2, printed pp.~34--35]{apple2024}. Together these qualifications
prevent a simple conversion of reported revenue into gross customer receipts.

Presentation creates a second distinction. Where Apple obtains control of a
third-party product before sale, it reports the gross amount billed. For
third-party applications sold through the App Store, its stated treatment
records only the retained commission in Services revenue. Services also already
includes recognition of the deferred value of services bundled with products
\citep[Note 2, printed p.~35]{apple2024}. Adding that recognition to reported
Services a second time would double count it; treating Services as gross App
Store customer spending would change the measured quantity.

The product and geographic tables describe different divisions of the same
sales. The five product/service rows sum to 391,035 million dollars in fiscal
2024 and 124,300 million in the quarter ended 28 December 2024. The geographic
segment sales sum to those same totals. They must not be added together.
Apple says the product mix is generally similar across segments, except that
iPhone represents a moderately higher share in Greater China. That qualitative
statement does not identify a product-by-region matrix
\citep[Notes 2 and 13, printed pp.~35--36 and 46--47]{apple2024}
\citep[Notes 2 and 10, printed pp.~6 and 11]{apple2025q1}.

\subsection{The movement that the note leaves unidentified}

At 28 September 2024 Apple reports total deferred revenue of 12.8 billion
dollars. At 28 December it reports 13.2 billion, and says that quarter's net
sales include 3.7 billion recognised from the deferred balance at the September
date \citep[Note 2, printed p.~6]{apple2025q1}. The balance sheet's December
current deferred-revenue line is 8,461 million, with 8,249 million at September
\citep[printed p.~3]{apple2025q1}. Those are current liabilities; the note's
larger amounts are totals. The precision differs too: the total is displayed
to tenths of a billion, whereas the current line is displayed in millions.

Let $D_0$ and $D_1$ denote opening and closing total deferred revenue. During
the quarter, let $A$ be additions to deferred revenue, $R_o$ revenue released
from the opening balance, and $R_n$ revenue released from amounts added during
the quarter. Let $M$ collect other signed movements, positive when they increase
the liability. These symbols describe a balance reconciliation; $M$ is not an
assertion that a particular adjustment occurred. In billions of dollars,
\begin{equation}
 D_1=D_0+A-R_o-R_n+M,
 \qquad A-R_n+M=D_1-D_0+R_o\simeq13.2-12.8+3.7=4.1.
 \label{eq:apple-deferred-movement}
\end{equation}
The source identifies $R_o$, not the whole release $R_o+R_n$. The 4.1 billion
residual therefore measures net additions after new-origin releases and other
movements at the displayed amounts. Calling it gross new deferrals, billings
or cash collected would be wrong relative to what has been calculated.

Rounding also matters. If each of the three displayed inputs is rounded to the
nearest 0.1 billion, a closed outer envelope for that residual is
$[3.95,4.25]$ billion: the lower endpoint uses $13.15-12.85+3.65$, and the
upper uses $13.25-12.75+3.75$. This is a conditional precision calculation,
not a verified issuer rounding rule or a bound on gross collections. Its
endpoints leave the treatment of exact rounding ties unspecified.

Two constructed histories make the missing movement visible. For this
illustration alone, use the displayed amounts as exact inputs and set $M=0$.
Keep the reported quarter's total revenue at 124.3 billion. The table supplies
two different choices for $A$ and $R_n$; other recognised revenue is the
remainder needed to match that total.

\begin{center}
\begin{tabularx}{\linewidth}{@{}Xrr@{}}
\toprule
Quarterly amount (USD billion) & History 1 & History 2\\\midrule
New deferrals, $A$ & 4.1 & 5.1\\
Release from those new deferrals, $R_n$ & 0 & 1.0\\
Release from the opening balance, $R_o$ & 3.7 & 3.7\\
Other recognised revenue & 120.6 & 119.6\\
Closing deferred revenue & 13.2 & 13.2\\
Total recognised revenue & 124.3 & 124.3\\\bottomrule
\end{tabularx}
\end{center}

Both histories satisfy the selected note totals, yet their new deferrals differ
by one billion. They are witnesses within this aggregate model, not reconstructed
Apple contracts or a fit to every financial-statement constraint. Additional
contract, cohort, return and movement information can rule out or refine them.
No probability attaches to either merely because two examples were written down.

The corresponding recognition entry debits deferred revenue and credits revenue
when service is provided. Receipt or billing of consideration that creates the
deferral instead credits deferred revenue and debits cash or, where an
unconditional right to payment exists, a receivable. Collection of that
receivable changes cash and receivables without creating revenue a second time.
For example, hold History 1's other revenue cash-settled and omit costs, taxes,
refunds and other movements solely to isolate this subsystem. Receiving its
4.1 billion of new consideration gives revenue-side cash inflow 124.7 billion.
If that new consideration is receivable at quarter end, cash inflow is 120.6
billion and receivables increase 4.1 billion. In both cases the asset increase
is 124.7, matched by a 0.4 increase in deferred revenue and 124.3 of revenue
before omitted expenses. These conditional entries explain the identification
problem; neither inflow is an estimate of Apple's actual cash collections.

\subsection{A dated recognition expectation is not a cash forecast}

At both September and December, Apple expects 64\% of the then-reported total
deferred revenue to be realised in less than a year, 25\% within one to two
years, 9\% within two to three years and 2\% after three years
\citep[Note 2, printed p.~36]{apple2024}
\citep[Note 2, printed p.~6]{apple2025q1}. These are two dated expectations
about two balances. Multiplication at the displayed December values gives
$13.2\times0.64=8.448$ billion in the first bucket. The corresponding September
calculation gives 8.192 billion. Their difference does not measure revenue
recognised during the intervening quarter: the dates, remaining horizons and
contract populations have changed.

The percentages divide an expected amount by recognition horizon. They do not
assign probabilities to alternative histories, fix collection dates or forecast
all future sales. The 8.448 billion calculation also cannot replace the reported
8.461 billion current liability. Rounded amount and percentage disclosures do
not justify an exact equality between those figures, and the query must retain
the source's distinct description of each quantity. A future calculation that
requires monthly service or receipts needs finer timing and payment evidence.

The note leaves another part of future service unquantified. For certain
long-term arrangements, Apple has not yet delivered the service, has no right
to bill for it and attributes the unbilled consideration entirely to that
undelivered service. It reports neither revenue nor an amount for those
obligations in this passage \citep[Note 2, printed p.~35]{apple2024}. A query
must retain an unknown amount; setting it to zero or equating it with recorded
deferred revenue would erase the distinction the company makes.

These disclosures therefore yield several different results. The total and
current deferred balances are reported observations. The 4.1 residual is a
derived net movement with unresolved components. The recognition buckets are
management's dated expectations, and the excluded consideration is unquantified.
Contract-level allocations, service lives, new-origin releases and collection
facts remain necessary for a particular future journal path. The annual and
interim periods also stay separate: the December 2024 quarter follows the
September 2024 fiscal year, so subtracting it from that annual revenue cannot
produce the year's final nine months. The complete revenue-note reading makes
these consequences explicit while the wider Apple disclosure package remains
only partly formalised.

\section{Apple: commitments and the cash needed to fulfil them}
\label{sec:apple-commitments}

The same distinction between recognition and cash arises when Apple buys
capacity and components. A lease liability measures one part of an existing
obligation; a purchase commitment describes goods or services still to be
acquired. A funding projection needs the payments arising from both, together
with new orders and operating decisions. Adding every disclosed amount would
mix measurements and could count the same requirement twice.

\subsection{From a lease liability to future payments}

At 28 September 2024, Apple's lease note reports the following amounts in
millions of dollars \citep[Note 8, printed p.~42]{apple2024}. The maturity table
first totals undiscounted payments and then deducts imputed interest to reach
the recognised liability. The right-of-use assets are reported separately.

\begin{center}
\begin{tabularx}{\linewidth}{@{}Xrrr@{}}
\toprule
Amount (USD million) & Operating & Finance & Total\\\midrule
Undiscounted lease payments & 13,518 & 1,272 & 14,790\\
Less: imputed interest & 1,984 & 376 & 2,360\\
Recognised lease liabilities & 11,534 & 896 & 12,430\\
Right-of-use assets & 10,234 & 1,069 & 11,303\\\bottomrule
\end{tabularx}
\end{center}

Thus $14{,}790-2{,}360=12{,}430$ reconciles the maturity and liability
measurements. The recognised liability is already included in the balance
sheet; it is not an extra cash requirement to add to 14,790. Its current
components total $1{,}488+144=1{,}632$ million, whereas the fiscal-2025 maturity
bucket is 1,991 million. A current carrying amount and a gross payment bucket
answer different questions. The right-of-use asset likewise need not equal the
remaining liability; these reported balances do not disclose the complete
amortisation, prepayment and modification history.

Apple also reports 849 million of fixed payments under additional leases that
had not commenced. It expects them to commence in 2025--2026, with terms from
less than one year to twenty-one years. Adding this separately disclosed amount
to the maturity total gives $14{,}790+849=15{,}639$ million. That sum is
consistent, under a nearest-0.1-billion rounding assumption, with the 15.6
billion of fixed lease obligations in management's cash-requirements discussion
\citep[printed pp.~25 and 42]{apple2024}. Numerical compatibility does not
establish a contract-by-contract reconciliation. It neither adds 849 to the
recognised liability nor identifies its commencement and payment dates.

The remaining lease information cannot supply those dates mechanically. The
note describes original terms typically no longer than ten years, with renewal
options, some reasonably certain to be exercised. Its weighted-average
remaining term of 10.3 years describes a different quantity. The reported
3.1\% discount rate, generally based on estimated incremental borrowing rates,
does not supply every contract's rate or the timing hidden in ``thereafter''.
Applying that one rate to invented year-end payments would introduce new
assumptions. The policy combining lease and nonlease components also has a
stated scope: corporate, data-centre and retail facilities
\citep[Notes 1 and 8, printed pp.~34 and 42]{apple2024}.

Fixed and variable costs add another distinction. For 2024, Apple reports
2.0 billion of fixed operating-lease cost, 1.9 billion of fixed operating-lease
cash payments and 1.0 billion of noncash right-of-use assets obtained for lease
liabilities. Variable costs across its leases were 13.8 billion; payments of
that kind primarily depend on purchases of output from the leased assets
\citep[Note 8, printed p.~42]{apple2024}. Historical expense is therefore not a
fixed future cash schedule. Forecasting it requires output, price and settlement
assumptions. Adding it to manufacturing commitments without identifying their
relationship could double count procurement requirements.

\subsection{Why a fall in commitments does not identify cash spent}

The purchase note and management discussion disclose different views. Note 12
lists 11,226 million of future payments under unconditional purchase obligations
whose remaining terms exceed one year. Its fiscal-2025 row is 3,206 million
\citep[Note 12, printed p.~45]{apple2024}. The term restriction selects
contracts; it does not place every payment beyond the coming year. The note
describes these as off-balance-sheet commitments, primarily involving supplier
arrangements, intellectual property, content and distribution rights.

Management's discussion provides additional near-term evidence. At the same
September date it reports manufacturing purchase obligations of 53.0 billion,
of which 52.9 billion is payable within twelve months. Other purchase
obligations total 12.0 billion, with 4.1 billion within twelve months; that
description includes capital assets as well as supplier, content and rights
arrangements \citep[printed p.~25]{apple2024}. A model confined to Note 12
would miss those reported aggregates. Conversely, subtracting 11.226 from 12.0
and labelling the difference short-term contracts is unsupported: the category
and term populations have not been reconciled, and the precision differs.

At 28 December 2024 manufacturing obligations are 39.8 billion, all payable
within twelve months. The interim report says contractual cash requirements
have not changed materially since the annual filing except for commercial
paper, manufacturing obligations and the State Aid tax payable
\citep[printed p.~16]{apple2025q1}. That qualitative statement is not an exact
December lease measurement or evidence that every lease remained unchanged.
The two manufacturing totals are dated observations over changing orders and
payment horizons. Their 13.2 billion difference does not identify purchases
or cash paid during the quarter.

To see what is missing, consider an explicitly simplified model in billions
of dollars. Let $O_0$ and $O_1$ denote a comparable population of undelivered
fixed-price orders at the two dates. Let $N$ be new orders and $F$ the value
removed when goods are delivered. Assume no cancellation, price or currency
change. Treat the displayed 53.0 and 39.8 as exact inputs for this illustration.
These assumptions define the example; the disclosure does not establish this
roll-forward convention for every reported obligation. The balance relation is
\begin{equation}
 O_1=O_0+N-F,\qquad F-N=53.0-39.8=13.2.
 \label{eq:apple-commitment-movement}
\end{equation}
With no new orders, delivery is 13.2. With 5.0 of new orders, delivery is 18.2.
Both leave the same closing undelivered amount of 39.8. The source balances
alone cannot choose between these conditional movements; actual contract
changes and delivery evidence could do so.

Payment is a further choice. Suppose the first history's 13.2 of delivered
goods all remain unsold inventory, with no opening payables, prepayments or tax.
Paying on delivery debits inventory and credits cash by 13.2: total assets and
profit do not change within this subsystem. Delivery on credit instead debits
inventory and credits trade payables by 13.2: assets and liabilities rise
together, with no immediate cash or profit effect. Settling that payable later
debits the payable and credits cash; consumption, sale or impairment would
require further entries. These examples isolate recognition and settlement.
They are not estimates of Apple's actual purchases or payments, or a fit to
all its financial statements. No probability follows from writing two histories.

Delivery assumptions also need the narrative that surrounds the purchase
table. Apple describes single or limited sources for some components,
industry-wide shortages and price fluctuations, and capacity or yield
constraints for new technologies. Suppliers may favour common components over
custom ones. Its named manufacturing locations describe outsourced production,
not a product-by-revenue-region allocation or ownership map
\citep[Note 12, printed p.~46]{apple2024}. A scenario in which input availability
changes may alter output, variable lease payments, inventory and receipts
together. The passage motivates those dependencies without identifying their
size or probabilities.

The same note, and the interim contingency note, preserve management's dated
assessment of unresolved legal claims. At each date, management states that
there was not at least a reasonable possibility that Apple had incurred a
material loss, or a material loss greater than a recorded accrual
\citep[Note 12, printed p.~46]{apple2024}
\citep[Note 9, printed p.~11]{apple2025q1}. It does not make litigation exposure
zero. Contract terms, supply constraints, commencement, delivery and settlement
must therefore remain in the historical record alongside reported balances.
They determine which future journal paths can be supported when a funding
question is asked.

\section{Apple: resources available for payment}
\label{sec:apple-liquidity}

The procurement and settlement choices just described determine when Apple
needs money. Funding those payments requires a second account: what resources
can be used, how securities become cash, and when customers pay. Annual and
interim Note 4 distinguish these questions even though a single total might
appear to answer them. The following table combines the notes' measurement
and balance-sheet classifications. Amounts are in millions of US dollars at
the displayed precision; the December date is the end of the first fiscal
quarter, not another annual endpoint
\citep[Note 4, printed pp.~36--37]{apple2024}
\citep[Note 4, printed pp.~7--8]{apple2025q1}.

\begin{center}
\begin{tabular}{lrrr}
\toprule
Resource / measurement &30 Sep 2023&28 Sep 2024&28 Dec 2024\\\midrule
Cash and cash equivalents &29,965&29,943&30,299\\
Current marketable securities &31,590&35,228&23,476\\
Non-current marketable securities &100,544&91,479&87,593\\\midrule
Total fair value &162,099&156,650&141,368\\
Total adjusted cost &173,752&160,600&147,088\\
Net unrealised gains / losses &$-11,653$&$-3,950$&$-5,720$\\\bottomrule
\end{tabular}
\end{center}

Cash plus the two securities rows equals total fair value. Adjusted cost plus
net unrealised gains and losses provides a second reconciliation of that same
total; it is not an additional resource to add to it. The source calls the
column \emph{adjusted cost}. Relabelling every investment as amortised-cost
debt would erase the different instruments represented in the table.

\subsection{Restriction changes what can be spent}

At September 2024, cash includes 2.6 billion dollars held in escrow, while
marketable securities include a further 13.2 billion. Both are restricted from
general use and designated for the State Aid tax obligation discussed earlier.
Subtracting their displayed centres from the total gives
$156{,}650-2{,}600-13{,}200=140{,}850$ million. MD\&A separately reports
unrestricted cash, equivalents and securities of 140.8 billion
\citep[printed pp.~25--26 and Note 4]{apple2024}. The apparent difference
does not establish conflicting disclosures: the restrictions and the MD\&A
total are rounded to tenths of a billion, whereas the table is in millions.
If each amount is rounded to its nearest displayed quantum, the subtraction
has possible bounds of 140,749.5 and 140,950.5 million. Those bounds overlap
the 140,750 to 140,850 million rounding range for the MD\&A total. This is a
conditional compatibility check, not an issuer-stated rounding convention
or a confidence interval.

The December footnote makes the distinction between maturity and restriction
especially clear: \emph{current} marketable securities include 3.3 billion
dollars of restricted escrow funds, and cash includes 126 million. At the
displayed centres, total resources less this disclosed escrow are
$141{,}368-3{,}300-126=137{,}942$ million. Cash plus current securities is
53,775 million; removing the same escrow gives 50,349 million. These are
indicative subtractions, not exact freely spendable balances. Under the same
nearest-rounding hypothesis the first subtraction ranges from 137,891 to
137,993 million. Interim MD\&A supplies no numerical unrestricted total
\citep[Note 4 and printed p.~16]{apple2025q1}. A current classification
therefore cannot serve as permission to spend an asset, and excluding escrow
does not establish immediate liquidation at the reported price.

Maturity supplies a different kind of timing information. Note 1 tests the
three-month cash-equivalent maturity at \emph{purchase}, while annual Note 4
classifies marketable debt as available-for-sale and assigns current or
non-current status using each instrument's underlying maturity
\citep[printed pp.~34 and 37]{apple2024}. Among non-current debt excluding
mortgage- and asset-backed securities, the September maturity shares are
86\%, 10\% and 4\% for 1--5, 5--10 and more than 10 years. The separate
mortgage- and asset-backed population has shares of 14\%, 9\% and 77\%.
At December these become 85/11/4\% and 13/10/77\%, respectively
\citep[Note 4, printed p.~8]{apple2025q1}. Each triple divides its own
population; the two triples cannot be added. Nor do contractual maturity
bands establish a sale date, sale price or cash receipt schedule. Management's
stated belief in sufficient resources also relies on ongoing operating cash
and continued access to debt markets. It is a dated funding assessment,
rather than a solvency guarantee derived from the table.

\subsection{Security turnover differs from valuation}

Securities can fund payments through sales or maturities, but a fall in their
reported balance need not equal cash released. Fiscal 2024 cash proceeds from
maturities and sales are 51,211 and 11,135 million, while purchases use 48,656.
Net cash released is therefore $51{,}211+11{,}135-48{,}656=13{,}690$ million.
The marketable-security balance falls by only
$126{,}707-132{,}134=-5{,}427$ million
\citep[cash-flow statement, printed p.~33, and Note 4]{apple2024}.
To expose the missing reconciliation, let $M_0$ and $M_1$ be total current
and non-current marketable securities at the period's opening and closing.
Let $C$ be maturity and sale proceeds less purchases, positive when cash is
released. Define $R$ as the unallocated difference between the stock movement
and that cash measure, all in millions over the same period:
\[
\begin{aligned}
 M_1-M_0&=-C+R,\\
 \text{Fiscal 2024:}\quad -5{,}427&=-13{,}690+8{,}263,\\
 \text{First quarter 2025:}\quad -15{,}638&=-13{,}335-2{,}303.
\end{aligned}
\]
For the quarter, the cash measure is
$15{,}967+3{,}492-6{,}124=13{,}335$ million
\citep[cash-flow statement and Note 4]{apple2025q1}. The relation defines a
residual; it does not explain it. Cash sale proceeds differ from the carrying
amount removed on disposal, and changes in measurement or classification may
also enter the stock. Apple's specific-identification policy for securities
sold requires the disposed lots to recover their cost
\citep[Note 1, printed p.~34]{apple2024}. Neither the positive
8,263 nor the negative 2,303 identifies a fair-value, currency or
reclassification journal.

The valuation columns help separate another tempting explanation. Across the
annual dates, total adjusted cost falls by 13,152 million and net unrealised
gains and losses rise by 7,703, producing a total fair-value decline of 5,449.
Across the quarter, the corresponding changes are $-13{,}512$ and
$-1{,}770$, producing $-15{,}282$. These totals include cash and equivalents,
whereas $M_0$ and $M_1$ include only marketable securities. They also differ from the
comprehensive-income statement's net-of-tax marketable-debt amounts: 6,054
million for the year and $-1{,}427$ for the quarter, including reclassification
adjustments \citep[comprehensive-income statement, printed p.~30]{apple2024}
\citep[comprehensive-income statement, printed p.~2]{apple2025q1}. Endpoint
unrealised balances and period OCI differ in population, tax and timing;
their numerical differences cannot be assigned to a tax charge or realised
gain without the missing movements.

Even reported cash needs matching populations. Annual cash-flow cash, cash
equivalents and restricted cash opens at 30,737 million, exceeding opening
balance-sheet cash and equivalents of 29,965 by 772. Both close at 29,943.
Thus the cash-flow change is $-794$, while the balance-sheet line changes by
$-22$; the cause and location of the earlier 772 remain unallocated. The
quarter's endpoint populations agree, giving the same 356 million increase
\citep[balance and cash-flow statements]{apple2024,apple2025q1}.

Valuation risk is likewise distinct from a payment. The notes describe Level 2
prices or models with significant observable inputs and counterparties that
\emph{generally} have high credit ratings. Annual Item 7A reports that a
hypothetical 100-basis-point rate rise at all tenors would reduce investment
fair value, including associated derivatives, by 2,755 million at September
2024 and 3,089 million at September 2023 \citep[printed p.~27]{apple2024}.
Those are dated issuer scenarios, not cash outflows or loss probabilities;
the issuer's valuation model has not been replicated. Interim Item 3's
no-material-change referral does not issue a new December sensitivity.

\subsection{Future receipts need customers and collection dates}

Receivables may turn into cash later, but their movement does not identify
collections. Annual net trade receivables rise by 3,902 million, while the
indirect cash-flow adjustment is $-3{,}788$. Their sum leaves an unallocated
114 million difference. In the quarter, the balance change of $-3{,}771$
and adjustment of $+3{,}597$ leave $-174$. Vendor non-trade changes match
the opposite indirect adjustments exactly: $+1{,}356/-1{,}356$ annually
and $-3{,}166/+3{,}166$ quarterly. Even that match supplies neither gross
sales nor total collections \citep[balance and cash-flow statements]{apple2024,apple2025q1}.
Sales, allowances, currency and noncash movements still need their own
evidence before receipts can be dated.

Concentration describes who might pay, but overlapping descriptions cannot be
summed as separate exposures. Carriers account for 41\%, 38\% and 33\% of
total trade receivables at the three successive dates. At December the note
also identifies two customers individually accounting for at least 10\%, with
shares of 18\% and 11\%. It does not say whether they are carriers
\citep[Note 4, printed p.~38]{apple2024}
\citep[Note 4, printed p.~9]{apple2025q1}. Let $A$ be those two customers,
$B$ the carrier group, and $s(S)$ the share of total trade receivables owed
by customers in a set $S$. Each customer belongs wholly to a set or does not.
Treat the displayed percentages as exact centres for this illustration.
Inclusion-exclusion gives
\[
 s(A\cup B)=s(A)+s(B)-s(A\cap B)
          =29\%+33\%-s(A\cap B).
\]
If both large customers are carriers, their 29\% is already inside the 33\%
carrier share, so the union is 33\%. If neither is a carrier, the union is
62\%. For a concrete construction, allocate receivables across eleven
hypothetical customers with shares
$18,11,9,9,8,7,4,9,9,8,8$ percent. These sum to 100 and only two reach 10\%.
Carriers can be the 18, 11 and 4 customers, totalling 33, or the first
9, 9, 8 and 7 customers, also totalling 33. The same marginal disclosures
therefore permit both endpoint unions. If only one large customer is a carrier,
the union is 44\% or 51\%; exact-centre membership does not permit an
arbitrary continuously divisible overlap. These are conditional percentage
examples, not recovered customer identities or probabilities of loss. The
note's trade denominator is not automatically the balance sheet's \emph{net}
receivable line, so multiplying it by that line would add an unsupported
measurement assumption. Required credit support or collateral applies to
\emph{certain} customers; it does not establish universal recovery.

Vendor non-trade receivables have a separate economic origin and denominator.
Apple buys components from suppliers and sells them to manufacturing vendors
that make subassemblies or final products for Apple. Those component sales
are excluded from product net sales; any gain reduces product cost of sales
when Apple sells the related final products. The two largest reported vendor
shares are 44\% and 23\% at September 2024, and 43\% and 24\% at December,
each pair totalling 67\% on its own dated vendor population
\citep[annual and interim Note 4]{apple2024,apple2025q1}. Equal combined shares
do not establish unchanged counterparties, component cohorts or collection
risk. A funding forecast consequently needs settlements and restrictions on
assets alongside payment dates on obligations. Keeping reported assertions
distinct from those missing facts is the starting point for representing
the company's historical record.

\section{Sales geography is not a map of costs or assets}
\label{sec:apple-segments}

The liquidity question leaves a related operating question. If demand weakens
in one place, which part of Apple's sales and profit is exposed? Apple reports
five geographic segments, but a segment is a management reporting division,
not a list of incorporated companies, factories or product lines. Reading its
two segment notes as though those were the same partitions would make a
regional scenario appear more precise than the filing permits.

\subsection{What the regional labels measure}

Apple's fiscal 2024 Note 13 names Americas, Europe, Greater China, Japan and
Rest of Asia Pacific. Americas covers North and South America. Europe also
includes India, the Middle East and Africa; Greater China includes China
mainland, Hong Kong and Taiwan. Apple says regional net sales are
\emph{generally} based on customer location and sales through retail stores
in those geographic locations. Its segments sell similar products and
services, but are managed separately around customers, distribution partners
and market conditions \citep[Note 13, printed pp.~46--47]{apple2024}. An
India-customer sales shock therefore belongs with Apple's Europe segment
under this disclosure's general attribution rule, even though India is in Asia. Neither this mapping
nor the label ``Greater China'' locates the legal seller, the manufacturing
partner or the cash-receiving entity.

The prior revenue reading established that product and regional net sales are
two classifications of the same 391,035 million dollars in fiscal 2024 and
124,300 million in the first quarter of fiscal 2025; the totals cannot be
added. The full Note 13 and interim Note 10 tables confirm five reported
regional sales and operating-income pairs in each dated column. For example,
Greater China has fiscal 2024 sales of 66,952 million and segment operating
income of 27,082 million; in the quarter ended 28 December 2024 the amounts
are 18,513 and 8,160 million
\citep[Note 13, printed p.~47]{apple2024}
\citep[Note 10, printed p.~11]{apple2025q1}. The quarter is a distinct
three-month period, not a fourth of the fiscal-year observation. Apple's
fiscal calendar can contain 53 weeks, as it did in fiscal 2023, and its
interim MD\&A describes historically higher first-quarter sales associated
in part with the holidays \citep[Note 1 and MD\&A]{apple2024,apple2025q1}.

Management attributes the annual Greater China sales decline mainly to lower
iPhone and iPad sales and reports an unfavourable renminbi effect. For the
first quarter of fiscal 2025 it attributes the decline mainly to iPhone,
while describing a favourable currency effect
\citep[annual and interim segment operating performance]{apple2024,apple2025q1}.
These are dated issuer explanations, not measured contributions of each driver
or an estimated response equation. A future region-specific demand or
exchange-rate path needs product, currency and contract information that the
regional totals alone do not supply.

\subsection{Why segment income cannot be treated as fully allocated profit}

Apple defines each segment's operating income as third-party sales less the
related cost of sales and directly attributable operating expenses. The
annual total of segment operating income is 162,044 million, but the
consolidated operating statement reports 123,216 million. Note 13 supplies
the bridge: 31,370 million of research and development and 7,458 million of
other corporate expenses, net, are subtracted. The latter category includes
general and administrative compensation, various nonrecurring charges and
other separately managed costs. The reported relation is
\begin{equation}
162{,}044-31{,}370-7{,}458=123{,}216
\quad\text{million dollars, fiscal 2024}.
\label{eq:apple-segment-income-bridge}
\end{equation}
The first-quarter fiscal 2025 bridge has the same form:
$53{,}174-8{,}268-2{,}074=42{,}832$ million. Both endpoints agree with the
consolidated statements of operations
\citep[statements of operations and Note 13, printed pp.~29 and 47]{apple2024}
\citep[statement of operations and Note 10, printed pp.~1 and 11]{apple2025q1}.
The annual ``other corporate'' amount is not identical to the 26,097 million
selling, general and administrative expense line in the consolidated
statement. Equating the two would change the reported classification.

Equation~\eqref{eq:apple-segment-income-bridge} identifies the company-wide cost pool, not how it would be
attributed to regions. To see the difference, take a hypothetical allocation
of the 38,828 million total corporate bridge across Americas, Europe, Greater
China, Japan and Rest of Asia Pacific of respectively 14,000, 10,000, 5,000,
4,000 and 5,828 million. Moving 1,000 from Americas to Greater China produces
a second allocation of 13,000, 10,000, 6,000, 4,000 and 5,828 million.
Both are nonnegative, sum to the same corporate total and reproduce the
consolidated 123,216 million. Yet Greater China's hypothetical profit after
these costs is either $27{,}082-5{,}000=22{,}082$ or
$27{,}082-6{,}000=21{,}082$ million. Apple does not report either fully
allocated figure, and the filing supplies no rule preferring one allocation.
These constructions establish numerical nonidentification within the stated
nonnegative illustration; they do not assert that all real cost allocations
are nonnegative or that either construction describes Apple's internal books.
Forecasting regional profit also requires the behaviour of directly
attributable and corporate costs under the scenario.

\subsection{Country sales and long-lived assets answer different questions}

The annual note separately presents country sales for fiscal 2024, 2023 and
2022 and country long-lived-asset stocks at the ends of fiscal 2024 and 2023.
It separately reports countries individually accounting for at least 10\% of
the respective totals and groups the rest; that selection rule does not make
``Other countries'' a single location. The 2024 displayed amounts illustrate
the different populations:
\begin{center}
\begin{tabular}{lrr}
\toprule
Country class & 2024 net sales & Long-lived assets\\
& \multicolumn{2}{c}{USD millions}\\\midrule
U.S. & 142,196 & 35,664\\
China, including Hong Kong and Taiwan & 66,952 & 4,797\\
Other countries & 181,887 & 5,219\\\midrule
Total & 391,035 & 45,680\\\bottomrule
\end{tabular}
\end{center}
The China country-sales row equals Greater China segment sales in all three
annual columns. The note's China footnote includes Hong Kong and Taiwan, so
the equality is a useful check on these displayed populations. It does not
turn that row into a legal-entity or production-location ledger. The U.S.
country-sales row is below Americas segment sales by 24,849 million in 2024;
assigning that difference to particular countries requires the precise
country-to-region attribution and customer facts. Europe's segment label
cannot be substituted for a European-country total because it includes India,
the Middle East and Africa \citep[Note 13, printed pp.~46--47]{apple2024}.

The 45,680 million long-lived-asset total equals net property, plant and
equipment at the 2024 balance date, and the 2023 totals both equal 43,715
million at their displayed precision
\citep[Notes 5 and 13, printed pp.~39 and 47]{apple2024}. The equality is
a checked comparison, not a definition of every asset's nature, owner or
productive capacity. Country asset stocks are measured at an instant, while
sales cover a year. Further, Apple says the information supplied to its chief
operating decision maker for segment assessment excludes asset information
\citep[Note 13, printed p.~46]{apple2024}. Thus neither a segment return on
assets nor an allocation of production capacity to the five sales regions is
identified by the disclosed tables. The statement that hardware is made
substantially by outsourcing partners located primarily in China mainland,
India, Japan, South Korea, Taiwan and Vietnam concerns another population,
and does not repair that missing asset/segment mapping
\citep[Note 12, printed p.~46]{apple2024}.

\subsection{The same totals permit different product exposure}

Note 2 reports 201,183 million of fiscal 2024 iPhone sales. Its four other
product and service rows sum to 189,852 million. Along with Greater China's
66,952 million regional total, these are \emph{margins} of an unreported
product-by-region table. Let $z$ be iPhone sales assigned to Greater China,
with all amounts in millions of dollars at the displayed centres. The four
cells for Greater China versus the other four regions must be
\begin{equation}
\begin{array}{c|cc}
 &\text{iPhone}&\text{Other products and services}\\\hline
\text{Greater China}&z&66{,}952-z\\
\text{Other four regions}&201{,}183-z&122{,}900+z
\end{array}
\label{eq:apple-product-region-margins}
\end{equation}
The row and column sums in~\eqref{eq:apple-product-region-margins} hold for
any nonnegative cells; they do not select
$z$. For example, $z=36{,}000$ and $z=38{,}000$ both meet the displayed
numeric margins. For an explicit extension, let $S_r$ be each other region's
reported sales; those four regions sum to $391{,}035-66{,}952=324{,}083$.
Assign each $(201{,}183-z)S_r/324{,}083$ of iPhone sales;
these amounts sum to the remaining iPhone total and are below $S_r$.
Then split every region's remaining sales across the four non-iPhone categories
in their disclosed company-wide proportions. This produces all five
product/service and five region margins exactly. These proportional splits
are mathematical witness constructions, not issuer allocations. The source adds that the product mix
was \emph{generally} consistent across segments, with a \emph{moderately}
higher iPhone share in Greater China. Both numerical examples put Greater
China's iPhone share above the combined other-region share. The qualitative
words have no disclosed tolerance, so neither matrix is certified as a
complete interpretation of that statement
\citep[Note 2, printed pp.~35--36]{apple2024}. The interim note repeats the
qualitative comparison for its own dated quarter
\citep[Note 2, printed p.~6]{apple2025q1}; it does not supply the missing
cross-classified cells.

If a scenario \emph{assumes} a 10\% decline in Greater China iPhone sales
alone, the two illustrative fiscal 2024 matrices imply 3,600 and 3,800
million less sales, respectively. The scenario percentage is a teaching
assumption, not an estimated shock; neither sales change determines profit,
cash flow, currency translation or future behaviour. The gap between those
two conditional answers comes from the unobserved joint allocation, even
before cost responses are considered. A reproducible model should retain
the reported marginals, the qualitative mix statement, any admitted
cross-classification and its source, and the unresolved allocation when
the evidence does not determine it.

Finally, the annual and first-quarter MD\&A state Apple's then-planned
fourth-quarter fiscal 2025 adoption of ASU 2023-07 with retrospective
transition \citep[Recent Accounting Pronouncements]{apple2024,apple2025q1}.
Those were dated plans, not the later disclosure itself; they do not
establish the standard's operative provisions without inspection of that
authority. The retained fiscal 2025 annual filing now supplies comparative
2024 segment expenses, examined in Section~\ref{sec:apple-pronouncements}.
A historical query keeps this later presentation separate from the older
filing's content. The regional calculations here are checked source-bound
comparisons and conditional witnesses, not a reconstructed Apple segment
ledger or a forecast-ready regional model.

\section{From geographic sales to legal entities}
\label{sec:apple-perimeter}

A regional sales forecast still leaves open which company sells the goods,
owns the inventory or owes the supplier. Those identities matter when a
scenario changes a particular entity's tax, funding or supply arrangements.
Apple's geographic segments cannot answer that question: their sales are
generally attributed by customer location and retail-store geography. The
consolidation policy and subsidiary exhibit supply a different kind of
information, about the companies represented in the group statements.

Annual Note 1 says that the consolidated financial statements include the
accounts of Apple Inc. and its wholly owned subsidiaries. The filing also
defines ``Company'' and ``Apple'' collectively in that way, unless otherwise
stated. Interim Note 1 repeats the wholly-owned scope for its condensed
statements and requires them to be read with the fiscal 2024 annual statements
and notes \citep[definition of Company and Note 1, printed pp.~1 and 34]{apple2024}
\citep[Note 1, printed p.~5]{apple2025q1}. This is positive evidence about the
reported consolidation perimeter. It is not a set of separate parent-company
accounts, nor does it identify every legal entity or ownership link within
that perimeter.

Exhibit 21.1 names the following subsidiaries and their jurisdictions of
incorporation. The labels below are transcribed from the exhibit, including
the separate U.S. state jurisdictions \citep[Exhibit 21.1, retained PDF p.~114]{apple2024}.
\begin{center}
\small
\begin{tabular}{ll}
\toprule
Subsidiary & Jurisdiction of incorporation\\\midrule
Apple Asia Limited & Hong Kong\\
Apple Asia LLC & Delaware, U.S.\\
Apple Canada Inc. & Canada\\
Apple Computer Trading (Shanghai) Co., Ltd. & China\\
Apple Distribution International Limited & Ireland\\
Apple India Private Limited & India\\
Apple Insurance Company, Inc. & Arizona, U.S.\\
Apple Japan, Inc. & Japan\\
Apple Korea Limited & South Korea\\
Apple Operations International Limited & Ireland\\
Apple Operations Limited & Ireland\\
Apple Operations Mexico, S.A. de C.V. & Mexico\\
Apple Pty Limited & Australia\\
Apple Sales International Limited & Ireland\\
Apple South Asia (Thailand) Limited & Thailand\\
Apple Vietnam Limited Liability Company & Vietnam\\
Braeburn Capital, Inc. & Nevada, U.S.\\
iTunes K.K. & Japan\\\bottomrule
\end{tabular}
\end{center}

The exhibit expressly limits the list. Its footnote states: ``Pursuant to
Item 601(b)(21)(ii) of Regulation S-K, the names of other subsidiaries of Apple
Inc. are omitted because, considered in the aggregate, they would not
constitute a significant subsidiary as of the end of the year covered by this
report.'' This preserves Apple's stated reason and date for the omission;
it does not establish a numerical significance test for this model. The
eighteen rows cannot be treated as an exhaustive group roster or as proof that
an omitted subsidiary is irrelevant to every entity-specific question.

The same filing makes that incompleteness concrete. Its tax-risk discussion
mentions Ireland and Singapore, ``where a number of the Company's subsidiaries
are organized,'' yet no Exhibit 21.1 row has Singapore as its jurisdiction
\citep[Risk Factors, printed p.~16]{apple2024}. The evidence supports
subsidiaries organized there while leaving their names and ownership paths
unresolved. Creating a company named ``Apple Singapore'' would turn a
geographic reference into an invented identity.

\subsection{Membership does not determine the ownership path}

Wholly-owned membership and immediate ownership answer different questions.
For a simple hypothetical example, a subsidiary could be owned directly by
the parent, or through another subsidiary that the parent wholly owns. If
every link is wholly owned, both arrangements preserve ultimate wholly-owned
membership while producing different immediate parents. This logical example
does not select either arrangement for any company in the table. Neither the
similar names of Apple Operations Limited and Apple Operations International
Limited nor their common incorporation jurisdiction establishes an ownership
link between them.

A model can therefore retain each named subsidiary, its reported jurisdiction
and the general consolidation assertion without drawing unsupported edges.
It should record an immediate parent, share class, ownership interest and
effective interval only when the evidence identifies them. The September 2024
exhibit is not a control history for earlier years or an updated list at
28 December 2024. Repetition of the general policy in the interim note does
not show that every subsidiary and ownership path remained unchanged.
Determining the dated group requires incorporation, acquisition, disposal,
reorganisation and control evidence as relevant; it cannot be recovered from
the fiscal calendar alone.

\subsection{Selling, manufacturing and owning are separate relationships}

Apple India Private Limited illustrates the distinction between incorporation
and sales attribution. The exhibit places its incorporation in India, while
Apple's segment definition places India in Europe. Under the general
customer-location rule, an India-customer sale belongs in that sales region;
neither statement identifies its legal seller or assigns all India sales to
that subsidiary. Even a complete subsidiary list would not supply the missing
seller-by-customer-region relationship.

The Business discussion also separates distribution channels: direct sales
through Apple's stores and sales force accounted for 38\% of fiscal 2024 net
sales, and indirect channels for 62\%. The latter include third-party cellular
carriers, wholesalers, retailers and resellers
\citep[Markets and Distribution, printed p.~2]{apple2024}. These are sales
channel shares, not ownership shares or counts of companies. A trading
relationship does not by itself establish that a counterparty belongs in the
consolidated group.

Production adds a further distinction. Apple says substantially all its
manufacturing is performed in whole or in part by outsourcing partners,
primarily in China mainland, India, Japan, South Korea, Taiwan and Vietnam.
It has also invested in manufacturing process equipment, much of which is
held at certain outsourcing partners, and made prepayments under long-term
supply agreements \citep[Risk Factors, printed p.~8]{apple2024}. Its MD\&A
describes partners manufacturing subassemblies and performing final assembly
and testing, alongside purchases from component suppliers
\citep[Manufacturing Purchase Obligations, printed p.~25]{apple2024}.
Equipment custody, a supply contract and corporate control are thus separate
relationships. A disruption at a partner may affect product supply, equipment
recoverability and prepayments without supplying evidence that the partner
should be consolidated. The affected assets, contracts and counterparties
must be identified before amounts can be assigned to that scenario.

\subsection{Tax and currency queries need facts beyond incorporation}

The jurisdiction column also leaves entity tax residence and functional
currency unresolved. Functional currency is the currency in which an entity
measures its accounts; it is a different question from the dollars in which
Apple presents its group statements or the currency of a particular invoice.
The group foreign-exchange discussion describes Apple as a net receiver of
currencies other than the U.S. dollar, with pricing and hedging responses to
exchange-rate movements. It supplies no entity-by-entity functional-currency
assignment \citep[Foreign Exchange Rate Risk, printed p.~27]{apple2024}.
An Ireland incorporation row therefore does not justify entering euros as the
entity's functional currency, just as the India row does not establish rupees.

Tax evidence has its own dates and units of analysis. Annual Note 7 describes
the historic tax allocation of profits of Irish branches of two subsidiaries
without naming those subsidiaries in that passage. The MD\&A describes the
resulting group payment obligation and restricted escrow
\citep[Note 7, printed pp.~39--40; State Aid Decision Tax Payable, printed p.~26]{apple2024}.
The four Ireland rows in Exhibit 21.1 do not identify which entities had those
branches, their historic tax residence or the jurisdictional tax bases used
in current accounts. A name match or country match would bypass the legal
and temporal facts needed for the tax calculation.

These distinctions determine what a supported query can return. The retained
sources answer questions about the stated consolidated scope and the named
subsidiaries' incorporation jurisdictions. Entity-specific tax, currency and
intercompany questions require additional records: legal counterparties,
dated ownership and control, compatible entity books, functional currencies,
tax bases and matched transactions. Group eliminations then require the
applicable accounting methods as well as those facts. The later IFRS group
examples explain their own assumptions; they do not establish the operative
US GAAP treatment for Apple. Keeping those dependencies explicit preserves
useful source knowledge while leaving the unanswered entity questions open.

\section{Legal proceedings beyond the booked liability}
\label{sec:apple-legal}

Knowing the companies within a group does not identify all the events that
could change their earnings. A legal proceeding can require a payment, but it
can also change how a company sells its products or earns licensing revenue.
Apple's reported App Store injunction illustrates the second channel: it
concerns developers' external purchase links on the U.S. storefront, rather
than a disclosed amount of damages. A forecast that represents legal exposure
only as a liability balance would leave that operating consequence outside
the model.

The annual Item 3 and interim Part II Item 1 each discuss Digital Markets Act
investigations, the Department of Justice lawsuit, Epic Games and other legal
proceedings. These are dated issuer reports, with different kinds of statements
inside each matter: allegations, reported decisions, pending motions,
management's position and possible consequences. Keeping those distinctions
allows a model to use what is reported without deciding the case itself
\citep[Item 3, printed p.~18]{apple2024}
\citep[Part II, Item 1, printed p.~18]{apple2025q1}.

\subsection{A loss assessment answers a narrower question than exposure}

Annual Note 12 and interim Note 9 give the same qualified assessment:
``In the opinion of management, there was not at least a reasonable possibility
the Company may have incurred a material loss, or a material loss greater than
a recorded accrual, concerning loss contingencies for asserted legal and other
claims'' \citep[Note 12, printed p.~46]{apple2024}
\citep[Note 9, printed p.~11]{apple2025q1}. The sentence identifies whose
assessment it is, concerns an incurred material loss, and allows for a recorded
accrual. It does not disclose a total legal accrual, allocate an accrual to one
of the named matters, or assert that all exposure is zero. Nor does it supply
a numerical probability or materiality threshold.

The annual MD\&A separately reports Apple's recognition policy: the company
records a liability when it is probable that a loss has been incurred and the
amount is reasonably estimable, a determination requiring significant judgment
\citep[Legal and Other Contingencies, printed p.~26]{apple2024}. That policy
explains why a pending proceeding and a recognised liability are different
facts. It is issuer-policy evidence; applying the operative US GAAP requirements
to a particular case still needs the relevant authority and case facts.
An allegation cannot supply those facts, and a possible change in future
sales is not itself the incurred loss described by the policy.

Settlement adds a further distinction. Annual Item 3 says certain matters
settled during the fourth quarter of fiscal 2024 without an individually or
aggregate material impact on financial condition or operating results. The
interim passage makes that statement about the first quarter of fiscal 2025.
These are two period-specific assertions, not the same settlement observation
copied forward. Neither passage identifies the matters, settlement amounts,
payment dates or earlier accruals. Both retain the warning that adverse
resolutions above management's expectations could materially affect the
reporting period. A payment schedule or an expense roll-forward therefore
remains unresolved even though the settlement statements themselves are known.

\subsection{The remedy, jurisdiction and procedural stage matter}

The Digital Markets Act discussion reports two investigations opened by the
European Commission on 25 March 2024. Article 5(4) concerns developers'
communication and promotion of offers and conclusion of contracts with end
users; Article 6(3) concerns default settings, app uninstallation and a browser
choice screen on iOS. The Commission's 24 June announcement is described as
preliminary findings alleging an Article 5(4) breach and the opening of a third
investigation into new contractual requirements for third-party developers
and app marketplaces. Apple states that it believes it complies and continues
to change its compliance plan in response to Commission engagement. The
preliminary allegation and Apple's position remain separate assertions.

Both filings say that, if the Commission makes a final violation determination,
it can issue a cease and desist order and may impose fines up to 10\% of
Apple's annual worldwide net sales. They also say an order would apply
immediately during an appeal unless a stay is granted. This preserves a
conditional consequence, its stated sales base and the stay qualification.
The percentage is not a fine already imposed, an expected loss or a booked
liability. Even a conditional fine calculation would need the applicable
operative law, relevant sales period and base, and assumptions about the
determination and remedy. None is supplied by multiplying a convenient annual
sales figure by the reported cap.

The Department of Justice discussion reports the civil antitrust lawsuit
filed on 21 March 2024 in the District of New Jersey by the DOJ and state and
district attorneys general. It alleges monopolization or attempted
monopolization in the markets for ``performance smartphones'' and
``smartphones.'' The DOJ seeks equitable relief; related state and federal
civil litigation seeks monetary damages and other nonmonetary relief. Apple
states that it has substantial defenses and intends to defend itself.
Treating all these proceedings as one quantified damages claim would erase
both the different requested relief and the absence of disclosed amounts.

Epic's case contains a reported court finding as well as unresolved procedure.
Apple says the Northern District of California court found certain App Store
Review Guidelines provisions to violate California's unfair competition law
and issued an injunction against prohibiting developers' external purchase
links. The stated scope is the U.S. storefront of the iOS and iPadOS App Store.
Apple reports implementing a compliance plan and filing its statement on
16 January 2024. Epic's motion disputing that plan and seeking enforcement is
reported as pending, opposed by Apple; Apple also reports filing a motion to
narrow or vacate the injunction on 30 September 2024. These passages do not
establish that the disputed plan was judicially accepted or quantify its sales
effect. An EU-wide scenario cannot inherit the U.S. injunction merely because
both matters concern the App Store.

The interim bodies repeat the substantive DMA, DOJ and Epic account, with
acronym or trademark presentation differences. Repetition is evidence of what
Apple reported in each filing, not proof that every intervening legal fact was
unchanged. The September 30 motion also falls after the annual September 28
period-end. Event date, reporting period and the date on which the report
became available must therefore remain distinct in a historical information
query.

\subsection{A partner's remedy can change Apple's revenue}

The risk discussion connects these matters to operations. Both filings report
changes to U.S. developer communications and EU changes to distribution,
payment processing, fee structures and developer tools. They also describe
licensing revenue from Google and other search providers, whose arrangements
are subject to investigations and proceedings
\citep[Risk Factors, printed pp.~14--15]{apple2024}
\citep[Part II, Item 1A, printed pp.~19--20]{apple2025q1}.
The interim adds that Google was found to have violated U.S. antitrust laws on
5 August 2024 and that the DOJ proposed remedies ranging from restrictions on
licensing agreements to prohibiting commercial search-distribution terms with
Apple for up to ten years. Apple says that, if implemented, those remedies
could materially affect its ability to earn that licensing revenue. The
finding is an earlier event newly described in this retained comparison;
the remedies are proposed, not reported here as imposed.

That channel differs from an Apple fine. The affected economic state would
include the search agreement, its revenue terms, duration, geography and
substitution possibilities. Likewise, an App Store remedy requires the
affected storefront, developer terms, fees, customer purchases and behavioural
response. The filings identify these channels without supplying the contract
economics needed to calculate them. They do not justify applying one loss
percentage to all Services revenue or treating every customer as affected.

A conditional calculation should consequently specify which matter and
remedy it assumes, when that remedy takes effect, which contracts it changes
and how the parties respond. A monetary-loss calculation additionally needs
recognition facts, any opening accrual and the settlement and payment schedule;
expense recognition and cash payment can occur in different periods. The
State Aid tax charge and escrow settlement, developed in
Section~\ref{sec:apple-tax-settlement}, retain their own tax records and must
not be counted again as a DMA loss. Until the missing
facts and applicable treatments are supplied, the model can retain the reported
proceedings and affected queries as unresolved, without assigning probabilities
or producing numerical journals. This preserves useful knowledge about future
exposure while keeping the distinction between a known assertion and a
supported financial consequence.

\section{When later disclosures change the historical query}
\label{sec:apple-pronouncements}

A later filing can reveal more about an earlier period without identifying a
different economic history. For example, a new segment-expense disclosure may
make a previously hidden cost visible, while the transactions that generated
that cost remain the same. To use that information correctly, a model must
retain the original report, the later presentation and the date at which the
additional detail became available. Apple's announcement and its later
segment tables show the difference between a planned disclosure and the
comparative information eventually reported.

Annual and interim MD\&A describe planned changes to tax and segment disclosures;
the interim adds expense disaggregation and a subsequent effective-date
clarification \citep[Recent Accounting Pronouncements, printed p.~26]{apple2024}
\citep[Recent Accounting Pronouncements, printed p.~17]{apple2025q1}.
The following table preserves Apple's stated adoption plans. The quarters are
fiscal, and the descriptions are issuer reports of the updates, rather than an
independent determination of their operative requirements.
\begin{center}
\small
\begin{tabular}{p{0.19\textwidth}p{0.40\textwidth}p{0.25\textwidth}}
\toprule
Update & Reported disclosure change & Apple's stated plan\\\midrule
ASU 2023-07 & Segment expenses and information about the chief operating
decision maker & Fiscal 2025 Q4, retrospective\\
ASU 2023-09 & Income-tax reconciliation and taxes-paid detail & Fiscal 2026 Q4,
prospective\\
ASU 2024-03, clarified by 2025-01 & Expense components and selling-expense
amount and definition & Fiscal 2028 Q4, prospective\\\bottomrule
\end{tabular}
\end{center}

The tax discussion says the rate reconciliation will include additional
information and detail for items meeting a quantitative threshold. Taxes paid
will be disaggregated into federal, state and foreign taxes, with further detail
for significant individual jurisdictions. It gives neither the threshold value
nor the future jurisdictional amounts. Those cash-tax populations also remain
distinct from tax expense, payable and deferred balances; their reconciliation
is developed in Section~\ref{sec:apple-tax-settlement}.
An announced disclosure change cannot be used to allocate an old tax total
among jurisdictions by assumption.

The earlier segment discussion concerns expenses that are significant and regularly
provided to the chief operating decision maker, together with that person's
title, position and use of segment profit or loss in assessing performance and
allocating resources. That anticipated disclosure could answer questions about
the information used by management. It did not yet identify the regional
allocation of corporate expenses or supply segment assets. The adoption
sentence was a dated plan at that information cutoff. The retained fiscal
2025 annual report now lets us check what followed.

Apple reports that it adopted ASU 2023-07 for its 2025 annual period using a
retrospective transition method \citep[Note 1]{apple2025annual}. Its Note 13
presents 2024 and 2023 comparative segment expenses. Take the Americas in
fiscal 2024. The earlier Note 13 reports net sales of $167{,}045$ million
and operating income of $67{,}656$ million, but no separate regional cost of
sales or selling-and-marketing amount \citep[Note 13, printed p.~47]{apple2024}.
The later comparative gives $89{,}587$ million of cost of sales and $9{,}802$
million of selling and marketing for that same region and fiscal year
\citep[Note 13, printed p.~47]{apple2025annual}. The displayed relationship is
\[
167{,}045-89{,}587-9{,}802=67{,}656\quad\text{(USD millions).}
\]
The 2025 table gives corresponding expense rows for the other four regions.
Together, their disclosed selling-and-marketing expenses total $18{,}639$
million. The five operating incomes still sum to $162{,}044$ million; after
the displayed corporate research and development of $31{,}370$ million and
general and administrative amount of $7{,}458$ million, the company reports
the same $123{,}216$ million of operating income. The 2024 report called its
$7{,}458$ million line ``other corporate expenses, net'' and qualified its
contents. The identical displayed amount is not, by itself, proof that the
two labels have identical scope. Both original and later assertions retain
their source identities.

This later presentation answers a prior-period disclosure question without
identifying new fiscal-2024 transactions or creating a correction journal.
Reading the named 2024 source leaves the regional expense detail unresolved;
reading the later source supplies its 2025 comparative assertions. An
investor-information view additionally needs evidence that the exact chosen
source version was publicly available by its requested cutoff. The accepted
corpus and filing-event metadata do not prove that for the retained 2024 PDF
or the later issuer-hosted HTML. Nor do the new regional expense rows reveal
the cross-classified product-by-region sales, entity cash allocation, or
behavioural response needed by the earlier regional stress example.

The interim expense passage reports ASU 2024-03's issuance in November 2024
and ASU 2025-01's clarification in January 2025. It describes purchases of
inventory, employee compensation, depreciation and intangible-asset amortization,
as applicable, within certain expense captions, a qualitative description of
the remaining amounts, and both the amount and the company's definition of
selling expenses. This describes disclosure populations; it does not provide
their values or make selling expenses identical to the entire selling, general
and administrative line. Purchases included in an expense disclosure likewise
cannot be identified with cash paid to suppliers without the relevant timing
and population facts.

The observed segment comparatives show why the two source-content views must
remain separate: the original cannot be overwritten even when a later filing
reports the same amounts with more detail. This comparison checks the affected
periods, dimensions, displayed units and issuer labels; it does not use the
phrase ``retrospective transition'' as an amount. The prospective tax and
expense plans still supply no backward allocation of missing components.
Their exact transition requirements need the operative updates and any
applicable amendments.

The January clarification illustrates the timing problem directly. It occurs
after the interim period ended on 28 December 2024 but is described in the
interim report. An event can therefore belong to the information in a report
without occurring inside its reporting period. A historical forecast cutoff
needs evidence of availability and the exact retained version, not merely the
period-end or the update's issuance month. Likewise, the earlier announced
segment adoption did not become an observed adoption until the later filing
was checked. The tax and expense plans remain future claims until their own
later evidence is inspected.

\subsection{Unchanged critical estimates do not mean constant amounts}

Annual Critical Accounting Estimates says management makes judgments,
assumptions and estimates affecting reported amounts, using historical
experience and other assumptions it considers reasonable to judge carrying
values. Its uncertain-tax discussion says reserves can be adjusted as facts
and circumstances change, and that final outcomes may differ from those
reserves. Its legal-contingency policy is the incurred-loss and estimability
condition examined above \citep[Critical Accounting Estimates, printed p.~26]{apple2024}.
These are descriptions of judgment and policy, not a disclosed numerical
estimation model or an assurance that the estimates will prove accurate.

The interim references both annual and interim Note 1 and reports no material
changes to critical accounting estimates since the annual filing
\citep[Critical Accounting Estimates, printed p.~17]{apple2025q1}.
That statement preserves the annual dependency and its materiality qualifier.
It cannot make every carrying amount, reserve or underlying assumption exactly
constant, and it supplies no numerical tolerance for a comparison. A model
should retain the reported policy and qualified continuity assertion alongside
the actual dated balances, then request the missing estimate or transition
facts when the query needs them. These estimate passages supply policy and
continuity assertions; they do not establish numerical stability, current
compliance or the tax and expense details still absent from these sources.

\section{Claims against another party can still affect Apple}
\label{sec:apple-litigation-risk}

The estimate policy leaves a further question: whose claim can give rise to a
cost for Apple? A company can face exposure through an agreement even when the
claim is against someone else. Apple's annual risk discussion says its
agreements sometimes include indemnification provisions that can subject it to
costs and damages in the event of a claim against an indemnified third party
\citep[Risk Factors, printed p.~13]{apple2024}. The direction matters: the
passage describes a possible obligation of Apple arising from another party's
claim, rather than compensation owed to Apple. A model limited to proceedings
naming Apple as defendant would miss that contractual channel.

The statement supplies neither an indemnified party nor a particular agreement.
It also leaves the covered conduct, trigger, duration and limit unknown. To
connect a third-party claim to a dated obligation, the model would need those
terms and the facts that activate them. The existence of indemnification
language somewhere in Apple's agreements cannot establish that a named case
is covered, that a payment is due, or that a liability has been recognised.
Those facts remain distinct from the management loss assessment examined in
Section~\ref{sec:apple-legal}.

The full annual passage also describes unresolved ordinary-course claims and
investigations, an increase over time in their number and alleged magnitude,
and patent claims concerning cellular-enabled products. Apple reports defending
infringement actions in U.S. and international courts, whose plaintiffs
frequently seek broad injunctions and substantial damages. These are issuer
statements about claims and requested relief. They do not identify a complete
matter population, decide the infringement allegations or quantify the
payments and product restrictions that will result.

Defense is a separate source of cost. The filing states, ``Regardless of the
merit of particular claims, defending against litigation or responding to
government investigations can be expensive, time-consuming and disruptive to
the Company's operations.'' Thus a scenario in which Apple successfully
resists a claim need not be a scenario with no defense effort or disruption.
The model therefore needs the response services, their amounts and dates,
and related payment obligations. The damages outcome cannot identify those
expenses or payments.

Settlement can alter recurring operations as well as resolve a dispute. Apple
says it may enter agreements or other arrangements to settle litigation,
without assurance of acceptable terms or that litigation will not occur.
Those agreements can increase cost of sales and operating expenses, change
business practices and limit products and services. A conditional settlement
therefore needs its payment terms and its continuing effects on royalties,
costs or offerings. An adverse resolution can additionally involve monetary
damages, disgorgement, corporate measures or an injunction. The annual passage
explicitly includes a resolution against an indemnified third party above
management's expectations. Neither that possibility nor the separately reported
quarterly settlements supplies a contract, amount or probability for a forecast.

\subsection{Insurance describes a possible recovery, not a net loss}

Apple says it maintains insurance for certain types of claims, while warning
that coverage may be insufficient for all losses or all claim types. The
qualification leaves both the covered population and the amount recoverable
unresolved. A model cannot deduct an assumed insurance payment from damages
or an accrued liability merely because insurance exists. It needs the policy
period, insured party and beneficiary, covered event, exclusions, deductible,
limit, insurer and recovery terms, together with facts about the claim. The
applicable recognition and presentation treatment must also be established
before a recovery amount becomes an accounting entry. The passages identify
none of those amounts or terms.

Other annual insurance statements concern different exposures. The
business-interruption discussion warns of recovery expenditures and lost sales,
especially when critical components have single or limited sources, while
saying insurance for certain types of losses may be insufficient
\citep[Risk Factors, printed p.~6]{apple2024}. The data-security discussion says
coverage is intended for certain aspects of those risks and may not cover all
losses or claim types \citep[Risk Factors, printed p.~12]{apple2024}. The trade
receivable discussion instead says a substantial majority of outstanding trade
receivables lack collateral, third-party bank support or financing arrangements,
or credit insurance \citep[Risk Factors, printed p.~16]{apple2024}. These
statements cannot be combined into one insured share of Apple's exposures.
The receivable qualifier is qualitative; it identifies neither every uncovered
customer nor a percentage. Insurance bought by Apple also does not make these
exposures an insurer's issued-contract measurement problem.

The interim risk preamble incorporates the annual risks and states that there
have been no material changes except those set out below
\citep[Part II, Item 1A, printed p.~19]{apple2025q1}. It preserves the annual
dependency and an express exception; it does not assert identical policies,
claims, coverage or loss amounts at the later date. Keeping the original risk,
the later qualified statement and the separately reported updates allows a
historical query to retain what was actually disclosed. The resulting account
connects direct proceedings, contract indemnities, defense, settlement,
operating restrictions and possible recovery without collapsing them into a
single loss balance. Calculating their consequences still requires the dated
contracts, claim facts and accounting treatment for each channel.

\section{A regulation can change the feasible operating plan}
\label{sec:apple-regulation}

A legal exposure need not begin with a claim or end with a payment. An import
restriction can prevent a product from reaching a customer, leaving the company
to change its supply route, offering or production plan. Apple's annual and
interim general-regulation risk both say that its products could be banned,
delayed or prohibited from importation despite procedures intended to satisfy
regulatory requirements. They also warn that measures may arrive with little
or no advance notice and that Apple may be unable to mitigate all adverse
impacts \citep[Risk Factors, printed pp.~13--14]{apple2024}
\citep[Part II, Item 1A, printed p.~19]{apple2025q1}. This describes a possible
constraint on operations, separate from the damages, defense and insurance
channels examined above.

The two reports preserve a broad list of regulated subjects: among them
antitrust, privacy, consumer protection, intellectual property, financial
services, trade, foreign exchange controls, cash repatriation, tax and
employment. The interim adds national security to that list. It also expands
the discussion of new and changing laws and regulations to name executive
orders, directives and enforcement priorities. These additions appear in the
cost-and-business-change, interpretation and specialized-application passages.
They identify the scope of Apple's later risk statement without identifying
a particular new instrument, effective date, jurisdiction or violation. A
model cannot create a dated legal event merely because a subject is newly
named in this comparison.

Both reports describe how changing requirements can increase costs, restrict
a product or feature, alter its design, affect customer demand and change the
business or supply chain. They describe further risks as products enter
specialized applications such as health and financial services or expand the
use of machine learning and artificial intelligence. The annual combines these
consequences and the compliance-policy discussion in one paragraph; the interim
splits and expands it. Paragraph count therefore cannot measure additional
economic exposure. Conversely, the notice-and-import paragraph repeats the
same substantive wording. That repetition establishes what Apple disclosed in
each report, rather than identical underlying laws, products or supply routes.

The policy qualification is essential to the mechanism. Apple says it has
implemented policies and procedures designed to ensure compliance, while giving
no assurance that employees, contractors or agents will not violate laws,
regulations or company procedures. Its separate import passage gives no
assurance that programs will prevent a violation or a claim of one. Policy
existence, compliance in a particular case, a claim of violation and effective
mitigation are thus different facts. Treating the reported procedures as proof
of compliance would remove precisely the uncertainty the source preserves.

A conditional operating scenario should start with the actual or assumed
measure and its effective date, legal scope, affected entity, product and
jurisdiction. It then needs the available supply routes, inventory, lead times,
customer demand and management response. For an assumed import delay, those
facts determine whether delivery is postponed, rerouted or lost, and which
storage, production or sourcing costs change. Recognition, payment and any
later claim must follow their own facts and treatments. The filings identify
these possible channels but supply neither a tariff rate nor a universal sales
reduction. Their reporting dates also do not establish when a particular rule
arose or when these exact report versions became available. Keeping disclosed
risk, operative authority, assumed restriction and response separate lets the
model retain a useful operating question without inventing its answer.

\section{Demand, disruption and recovery follow different paths}
\label{sec:apple-disruption}

An open supply route is useful only if a customer still wants the product.
Conversely, a customer may be ready to buy while a supplier cannot deliver a
critical component. The operating scenario introduced above must therefore
describe both demand and the ability to serve it. Apple's annual risk discussion
connects global and regional economic conditions to consumer spending, the
financial condition of its business partners and its own access to finance.
Its separate interruption discussion describes obstacles to manufacture,
delivery and service provision \citep[Risk Factors, printed pp.~5--6]{apple2024}.
These passages give reasons to retain several possible operating histories
behind the same reported financial totals.

Apple says that sales outside the United States account for a majority of net
sales and that a majority of its supplier facilities are outside the United
States. The two statements concern different populations. A sales region does
not locate a factory, and a factory's location does not identify the customers
whose orders depend on it. The segment definitions examined earlier generally
attribute sales by customer and retail-store location, while manufacturing
purchase obligations describe purchases from outsourcing partners and component
suppliers \citep[Note 13, printed p.~46; Liquidity and Capital Resources,
printed p.~25]{apple2024}. Linking the two requires the actual product,
component, supplier and route relationships. Geographic labels alone cannot
create them.

The macroeconomic passage describes several possible effects of the same
economic conditions. Recession, unemployment, inflation, tighter credit,
higher interest rates and currency fluctuations can weaken consumer confidence
and spending. Fiscal and monetary policy, market volatility and changes in
income or asset values can also affect demand. Suppliers, manufacturers,
logistics providers, distributors, carriers, other channel partners and
developers may instead face financial instability, difficulty obtaining credit
or insolvency. Apple further identifies trade-receivable collectibility,
derivative-counterparty and financial-institution failure, limited access to
new debt, reduced liquidity and declining financial-instrument fair values.
The disclosure permits a scenario to connect these channels through common
conditions; it supplies neither a joint probability law nor a requirement
that every channel deteriorate together.

An interruption raises a different operating question: which activity becomes
unavailable, and for how long? Apple's discussion includes political and trade
disputes, geopolitical conflict, terrorism, natural disasters, public-health
issues and industrial accidents. It says substantially all manufacturing is
performed wholly or partly by outsourcing partners located primarily in China
mainland, India, Japan, South Korea, Taiwan and Vietnam. Trade restrictions can
raise input costs, require different suppliers or business relationships, or
lead the company to cease offering affected products, services or third-party
applications. Changing the supply chain may itself be costly and disruptive,
and restrictions may arrive with little notice. Apple reports that past
U.S.--China tensions led to tariffs and other restrictions affecting its
business. This historical assertion is distinct from its conditional warning
about further escalation; a future scenario still needs a specified measure,
effective date and affected trade.

The same passage describes facilities in disaster-prone locations and asserts
that climate change is increasing the frequency or intensity of certain
natural disasters and extreme weather. It also identifies fire, power
shortages, nuclear and other industrial accidents, hostile acts, ransomware,
labor disputes and health emergencies. These are Apple's stated exposures,
without facility-specific hazard estimates. Its requirement that suppliers
maintain safe workplaces is separate from whether an accident occurs. The
company reports that major public-health issues, including COVID-19, have
adversely affected demand, operations and distribution, and warns of possible
future effects. Travel restrictions, freight limitations and interruptions to
production ramps supply concrete mechanisms; the paragraph does not quantify
an historical loss or a future duration.

Consider two explicitly assumed histories for a particular product. In one,
the supply route remains available but customers reduce orders. In another,
orders remain but a critical component cannot arrive. The first requires a
demand response; the second requires an inventory and delivery response.
They can also occur together. To compare the histories, the model needs
product-level orders, available components and finished goods, lead times,
contractual commitments and feasible management actions. An assumed alternate
supplier must have capacity and suitable terms before it becomes an available
response. A sales delay or lost order follows from those facts and assumptions;
it cannot be obtained by applying one percentage to every reported balance.

Recovery adds another date. Apple warns that an interruption can require
substantial time and expenditure to resume operations, with lost sales and
greater consequences where critical components have single or limited sources.
Note 12 separately explains that even components available from several
suppliers can face industry-wide shortages and commodity-price fluctuations
\citep[Note 12, printed p.~46]{apple2024}. Ending an assumed interruption
therefore need not restore production or sales immediately. Recovery spending
can precede or follow resumed deliveries. The model must retain the dates of
the interruption, restored capacity, deliveries, expenditure and settlement
separately; the applicable accounting treatment then determines the journals.
The insurance qualification already examined in the litigation discussion
does not collapse those dates into a net loss or guarantee reimbursement.

The reported receivables, investments, debt and liquidity resources constrain
particular dated histories, but they do not determine these operating
responses. The investment sensitivity discussed above measures a specified
interest-rate effect. Chapter~\ref{ch:financing-groups} examines the distinct
debt and currency measures; none identifies the broader disruption response.
The interim's incorporation of annual risk factors, except for the changes it
sets out, likewise does not establish unchanged facilities, orders or funding
conditions \citep[Part II, Item 1A, printed p.~19]{apple2025q1}.
For the historical record, the useful distinction is now precise: retain
Apple's assertions, the relationships needed to interpret them, and the
analyst's scenario assumptions as separate information. Alternative demand,
supply and recovery histories can then remain visible until evidence narrows
them.

\section{Product transitions change the relevant history}
\label{sec:apple-competition-transition}

Keeping a supply route open does not ensure that it supplies the product
customers now want. A new offering can replace demand for an earlier one while
the company still holds its components, finished goods or purchase commitments.
Apple's product-transition risk explicitly says that new products, services
and technologies may replace existing offerings and produce lower revenues
and lower profit margins \citep[Risk Factors, printed p.~7]{apple2024}.
A model that equates a successful release with an improvement in financial
performance would remove a possibility the source asks it to preserve.

The preceding competitive-market discussion explains why a transition is not
merely a production date. Apple describes price competition, short product
life cycles, changing standards, rapid adoption of technology and price-sensitive
customers. It says that developing almost the whole product solution requires
substantial research and development, without assurance of the expected returns
or successful marketing. Its reported intellectual-property portfolio also
does not assure protection in every jurisdiction. The discussion of competitors
imitating products or infringing rights remains an issuer assertion; it does
not decide a particular infringement dispute. Apple's description of minority
share in three global device markets and of competitors able to sell at little
or no profit supplies competitive context, rather than a numerical market
share or a forecast of Apple's margins. Services face a further set of
alternatives, including free content. The company's expectation of more
intense competition is a dated expectation, not an observed future result.

The transition passage identifies several separate conditions for a new
offering: development must succeed on time, customers must accept it, production
must reach suitable quantities and costs, and application software or other
third-party support must be available. New technology can create ramp-up or
quality problems. Purchase commitments and inventory must also be managed in
line with anticipated demand. Note 12 makes the supply condition concrete:
new products often use custom components available from only one source;
capacity can be constrained while yields mature or manufacturing expands,
and a supplier may instead concentrate on common components
\citep[Note 12, printed p.~46]{apple2024}. A forecast of customer acceptance
cannot establish that those components will be available.

Suppose, as an illustration, that customers switch from an earlier product to
its replacement. If existing components can be used in the replacement and
the supplier can expand output, the transition presents one set of feasible
production choices. If the components are incompatible or the new component
has constrained capacity, the same assumed customer response presents another.
Neither history is identified by aggregate revenue or inventory. The comparison
requires the old and new product definitions, component compatibility, inventory
cohorts, committed purchases, cancellation terms, lead times and available
capacity. Prices, quantities and costs would then be separate scenario inputs
or model consequences. The risk language gives no substitution fraction,
write-down amount or probability with which to combine the histories.

Even a delivered replacement has more than one financial date. Apple's revenue
policy generally connects product revenue to transfer of control, while
allocated revenue for bundled services is deferred over the estimated service
period; collection follows its own terms
\citep[Note 2, printed p.~35]{apple2024}. Development, announcement, production,
shipment, service provision and cash receipt therefore cannot share a single
``launch'' timestamp. The reported policy must be applied to the relevant
contract and performance obligations. Any inventory disposition, commitment
change or development-cost treatment needs its own facts and applicable
accounting method. Preserving these distinctions lets the model ask whether
an assumed transition is feasible and how its financial consequences would
arise, while leaving actual product histories unresolved.

\section{A supply agreement does not establish available output}
\label{sec:apple-component-availability}

Once a transition identifies the components a product needs, the next question
is whether they can arrive on the required date. An agreement to purchase them
does not answer that operating question. Apple's component-availability risk
says that supply agreements cover many components but gives no assurance that
they can be extended or renewed on similar terms, or at all. It also identifies
supplier financial difficulties, business failure and industry consolidation as
possible limits on supply \citep[Risk Factors, printed p.~8]{apple2024}.
The historical record must therefore distinguish the agreement, its effective
term and the supplier's ability to perform. Renewal and delivery need their
own evidence before a scenario can rely on them.

More suppliers do not eliminate every shortage. Apple describes significant
supply and pricing risks for single or limited sources, while warning that
even components available from several sources can face industry-wide shortages
and commodity-price fluctuations. Its example of past semiconductor demand
and supply shortages reports an adverse effect on obtaining components and
products on commercially reasonable terms. The warning that such disruptions
could recur preserves a future possibility; the passage supplies no duration,
quantity shortfall or loss with which to calibrate it. A scenario may specify a
common industry constraint, but it must state that assumption rather than
derive independence or a diversification benefit from the number of suppliers.

Custom components present a different constraint. A sole source may initially
have insufficient capacity while yields mature or manufacturing expands.
Alternatively, the supplier may give priority to common components. A yield
problem and a production-allocation decision can both delay the same product,
yet the possible responses differ. Apple separately reports past adverse
effects from delayed or constrained component supply and delayed shipments of
completed products by outsourcing partners. Component production, final
assembly and delivery remain separate activities; an improvement in one need
not establish completion of the others. The Note 12 supply discussion and the
outsourcing passage retain the same distinction between sources, capacity,
partner performance and control \citep[Note 12, printed p.~46; Risk Factors,
printed p.~8]{apple2024}.

For an explicitly assumed delay, compare restoring the current source with
switching to another. The disclosure says that consequences can depend on the
time needed to obtain sufficient quantities from either route. A named
alternative is therefore only a proposed response until component
compatibility, available quantities, lead times, price and contractual terms
are established. Existing purchase obligations do not identify those facts or
the release of the old commitment. The model should retain both responses
with their separate prerequisites and delivery dates, leaving an infeasible
or unresolved response visible. Only after the relevant events and accounting
facts are specified can the resulting purchases, settlements or inventory
movements generate journals. The filing does not select a route or assign
probabilities to the alternatives.

\section{Available output can still create losses and obligations}
\label{sec:apple-quality-inventory}

Obtaining components and completing delivery do not establish that a product
can safely serve its intended use. Apple's defects discussion covers its
hardware and software, third-party components and products, and service
outages, slowdowns and errors. It reports that services have sometimes failed
to perform as anticipated. It also warns that defects can create personal-
injury, property or environmental risks, that specialized uses such as health
can increase those risks, and that complex technologies such as artificial
intelligence can expose users to harmful or inaccurate content. Detection and
repair of every defect are not assured \citep[Risk Factors, printed p.~9]{apple2024}.
For the model, a delivered unit, an acceptable unit and a repaired unit are
different states. The source identifies possible transitions between them,
not a defect rate or a particular incident.

The financial consequences can extend well beyond inventory. Apple names
product liability, recalls, replacement or modification, write-offs of
inventory, property, plant and equipment or intangible assets, warranty and
other expenses, litigation costs and regulatory fines. Quality can also affect
reputation, acceptance, demand and the timing of new offerings. These effects
must not all become one inventory expense. A customer obligation and a possible
supplier recovery are also separate. The outsourcing passage says partner
arrangements may include defect-expense reimbursement while Apple generally
remains responsible to the consumer for warranty and out-of-warranty service
\citep[Risk Factors, printed p.~8]{apple2024}. Identifying a supplier does not
establish a reimbursement right, amount or receipt date.

Excess stock raises a different question even when its quality is adequate.
The inventory-risk passage says Apple records write-downs for product and
component inventories that are obsolete, exceed anticipated demand, or have
cost above net realizable value. It separately describes necessary cancellation
fee reserves for orders of excess products and components. Inventory already
held and an order not yet received are different exposures; cancellation
depends on the actual order terms. Note 1's first-in, first-out policy identifies
the reported cost-flow method, rather than a lot's realizable value or a
cancellation charge \citep[Risk Factors, printed p.~9; Note 1, printed p.~34]{apple2024}.

Apple also describes impairment reviews when circumstances indicate that
long-lived assets, including capital assets held at suppliers and inventory
prepayments, may not be recoverable. It reports a write-down for excess carrying
value over fair value if impairment is determined. This is the issuer's
description of its assessment; an executable calculation still needs the
actual asset classification, assessment basis and applicable accounting method.
The company's belief that the named assets and commitments are currently
recoverable is accompanied by no assurance against future charges. It neither
proves a future impairment nor permits the model to erase that possibility.
An inventory cost-flow calculation, a cancellation obligation and an asset
impairment assessment cannot be substituted for one another.

The link to product transitions is the forecast used before delivery. Apple
orders components and builds inventory ahead of announcements and shipments;
manufacturing purchase obligations cover forecast requirements, typically for
periods up to 150 days. It warns of excess or insufficient output, or failure
to use firm commitments fully, when forecasts are wrong. The qualifier
``typically'' prevents that period from becoming a universal contract limit.
Nor does the description make the separately reported Note 12 obligations
with remaining terms above one year the same population as manufacturing
obligations \citep[Risk Factors, printed p.~9; Note 12, printed p.~45]{apple2024}.

Consider an explicitly assumed reduction in expected demand. Two companies
could hold the same reported inventory amount while one holds saleable
components for the new offering and the other holds incompatible earlier
components, defective finished goods or commitments it cannot cancel cheaply.
Their feasible responses and possible losses differ. Lot identity and cost,
quality, compatibility, forecast dates, anticipated demand, realizable value,
order terms and customer or supplier obligations distinguish the histories.
Recognition and measurement then require their own evidence. The disclosure
does not supply those facts or probabilities, so the model preserves the
conditional histories without inventing a reserve or treating the forecast as
observed customer demand.

\subsection{Service provision requires continuing access}
\label{sec:apple-service-access}

Adequate components and saleable inventory still leave a service question:
can the company continue offering the software and content its customers want?
Apple says its products and services require third-party intellectual property,
while giving no assurance that the necessary licences will be available on
commercially reasonable terms or at all. Its digital-content arrangements can
be short and need not renew on similar terms. Content providers can also require
new rights-management or security technology before access can continue
\citep[Risk Factors, printed pp.~9--10]{apple2024}. The relevant resource is
therefore partly a dated right to offer particular content or functionality.
Physical capacity alone cannot establish that right.

Developer support has a separate economic cause. Apple believes customers'
hardware purchases depend partly on available applications, whose continued
development and maintenance are not assured. It describes developers comparing
platforms by market position, anticipated revenue, expected growth and the cost
of developing and updating software \citep[Risk Factors, printed p.~10]{apple2024}.
An application can consequently become unavailable or unsuitable even while
the hardware remains usable. A scenario needs the affected applications and
developers' response; an estimate of component capacity cannot supply either.

Access also depends on how a service uses data and takes payment. Apple reports
restrictions on collecting, using and transferring personal information,
including transfers within its international group. It separately warns that
failure to follow payment-card security standards could increase transaction
costs or remove the ability to process cards
\citep[Risk Factors, printed p.~15]{apple2024}. These issuer descriptions
identify possible constraints. Which constraint applies depends on the
data, entity, jurisdiction, agreement and effective rule.

Suppose a content agreement permits a specified offering only until its expiry
date. Under one explicitly assumed history it renews; under another it expires
without replacement. Identical physical resources then permit different
offerings. The comparison needs the service, customer population, territory,
contract term and allowed substitutions, together with any new cost and demand
response. A financial calculation must apply the relevant revenue, payment
and remedy terms to the service actually provided. Apple's risk statements
supply neither those contract facts nor probabilities for the alternatives.
The historical record must preserve the continuing relationships as well as
the reported balances if it is to support these different futures.


\section{A useful answer before the company's history is complete}
\label{sec:real-package-connection}

The Apple examples now give us enough information to pose a more demanding
question: what can one connected answer say about the reported company and
its possible cash requirements? The answer must carry both the reported
amounts and the limits of the histories used to explain them. Otherwise a
complete-looking statement could conceal a missing liability, a hypothetical
retirement or an unavailable settlement date.

For the retained annual source, the selected statements report net income of
93,736 and other comprehensive income of 4,280, giving comprehensive income
of 98,016 million dollars. Assets of 364,980 less liabilities of 308,030 give
closing equity of 56,950. The following quarter reports net income of 36,330,
other comprehensive income of 383 and closing equity of 66,758
\citep[consolidated financial statements]{apple2024,apple2025q1}. These are
source-located readings of displayed amounts. Their units, periods, group
perimeter and statement populations accompany them. A share count or a
per-share amount keeps its own unit; it cannot enter a monetary total merely
because it sits next to one.

The connected result retains the 102 selected annual rows, 100 quarterly rows
and their separate opening balance-sheet comparisons. It also retains all
23 note dispositions, 70 known references, seven package obligations,
32 proposed source assertions and seven additional frontier entries from
the frozen inventory. A note's remaining facts stay attached to it. These
records preserve the declared package; they neither discover every possible
reference nor establish that every retained sentence has a formal meaning.

A query for the selected displayed statement tokens can therefore finish even
though the question about actual whole-company histories remains unresolved.
The family question specifies the issuer, sources and target histories, but
has no admitted member list or probability measure. The distinction matters:
an unresolved feasible remainder does not mean that no history exists, and
two illustrated histories do not mean that the real company has exactly two
possible histories. The combined answer is partial because the financial
explanation has obligations that the reading task does not.

Preserving a source token also requires restraint in arithmetic. The annual
cash-flow row for debt issuance displays a dash. Its observation remains a
reported dash. The retained financing reconciliation separately states the
narrow convention under which that dash contributes zero to its displayed
sum. This convention does not turn all dashes into zeros. Similarly, the
annual cash-flow opening amount of 30,737 differs from the opening balance-sheet
cash amount of 29,965 by 772. That population difference is an explanatory
residual, not permission to create a receipt or payment for 772.

\subsection{Matching annual rows leaves different liabilities to settle}

Two conditional asset histories illustrate why the connected result must
retain journals as well as totals. In H1, gross additions of 10,529 are paid
with 9,447 of cash and a capital payable of 1,082. Retirements remove gross
cost of 6,000 and accumulated depreciation of 5,636, leaving a loss of 364.
H2 adds and retires another 1,000 of gross cost: additions become 11,529,
retirements 7,000, loss 1,364 and capital payable 2,082. Depreciation is 8,200
in both branches. The extra acquisition and retirement cancel in net property,
plant and equipment; cash is unchanged. The asset journals alone leave H2
with 1,000 more liability and 1,000 less profit and equity.

Now suppose a distinct non-tax expense liability exists: its amount is $q$
in H1 and $q-1{,}000$ in H2, with $q\geq1{,}000$. Each branch debits expense
and credits that liability. This is an independent economic premise, not an
event inferred from a statement residual. Suppose also that the retirement
loss and this expense share the product-cost row, both liabilities share the
same other-current-liability bucket, the extra asset cost has the same gross
class, and current and deferred tax, all other activities, shares and
distributions are unchanged. Then the expense difference of $-1{,}000$
offsets the extra retirement loss, and the liability difference offsets the
extra capital payable.

The indirect cash-flow presentation still needs a specified allocation. Let
$\alpha\in[0,1]$ allocate the expense-liability adjustment once between the
noncash adjustment and the working-capital adjustment. If $\Delta$ means
H2 minus H1, the two displayed-row differences become
\begin{align}
\Delta\text{other adjustments}
 &=1{,}000+\alpha(-1{,}000)=1{,}000(1-\alpha),\\
\Delta\text{other liabilities}
 &=(1-\alpha)(-1{,}000)=-1{,}000(1-\alpha).
\label{eq:real-package-cash-presentation}
\end{align}
The two differences sum to zero for every allowed $\alpha$. At $\alpha=1$,
each is zero and all 102 selected annual row differences vanish under these
assumptions. At $\alpha=0$ or $1/2$, the same total operating cash has a
different presentation in those two rows. Here $\alpha$ is a reporting
assumption, not a probability over histories. Matching differences also does
not construct an absolute feasible remainder for the company. The illustrative
journals remain conditional subsystem histories.

Their future consequences remain joint. With $q=2{,}000$, carrying-value
settlement of the two liabilities gives the following changes, in millions,
if both payments occur in the first future period:
\begin{center}
\begin{tabular}{lrrr}
\toprule Branch & Operating cash & Investing cash & Total cash\\\midrule
H1 & $-2{,}000$ & $-1{,}082$ & $-3{,}082$\\
H2 & $-1{,}000$ & $-2{,}082$ & $-3{,}082$\\\bottomrule
\end{tabular}
\end{center}
Paying the capital liability is an investing movement; paying the separately
accrued expense liability is an operating movement. Each payment debits its
liability and credits cash. Under the explicit assumption of no incremental
tax, currency or interest effects, it creates no new profit, other
comprehensive income or equity movement. Sufficient unrestricted funding is
assumed; no opening cash balance is manufactured.

Taking the lowest operating amount from H1 and the lowest investing amount
from H2 would produce $-4{,}082$. Neither history produces that total. The
result must return the two pairs, not independently selectable marginal
endpoints. Moving capital payment to the second period changes first-period
cash to $-2{,}000$ in H1 and $-1{,}000$ in H2; moving the expense payment
instead changes it to $-1{,}082$ and $-2{,}082$. Thus equal eventual cash does
not identify its timing. The six specified paths retain both branch identity
and payment period.

When dates are missing, the response keeps the conditional eventual amounts
and marks period timing unresolved. It distinguishes these calculated amounts
from reported tokens. A request for future depreciation remains unsupported
without the asset cohorts, tax state and other required facts. A probability
request remains unsupported without an admitted law. Neither refusal erases
the amounts that a narrower question can answer.

\subsection{The information used belongs to a particular view}

The later annual source provides a different kind of progress. Its FY2024
Americas comparative shows sales of 167,045, cost of sales of $-89,587$ and
selling and marketing expense of $-9,802$, giving operating income of 67,656
\citep[Note 13, FY2024 Americas comparative]{apple2025annual}.
The selected original source lacks those two separate expense rows. An
original-source query therefore returns unresolved components, while a query
explicitly naming the later source can return them. The result links the later
view separately; it does not silently enlarge the original view or assert
that measurement scope was unchanged. Later detail is not, by itself, an
accounting-error correction.

Named-source access answers what retained bytes contain. Investor-information
and system-as-of queries additionally require accepted clocks for those exact
bytes. Where publication, ingestion or review time is unestablished, the
historical query remains partial with a named timing obligation. A filing
event alone cannot date a separately converted PDF byte for byte. This prevents
the later comparative detail from becoming evidence that an earlier investor
could already have used.

JPMorgan's reference conflict illustrates the same discipline in another
company. In the retained 2023 report, Item 15 faithfully says page 156; Item 8
and the auditor heading support printed page 163, physical page 165, for
navigation \citep[Items 8 and 15; printed pp.~156 and 163]{jpmorgan2023tenk}.
The wrong target is printed page 156, physical page 158. Both use
the same two-page offset, so changing the offset cannot reconcile the targets.
A proposed navigation correction preserves the literal statement. If a
counterfactual earlier view had accepted the wrong target, its selected page
set, bounded search and coverage summary would need reconsideration, while
the literal and an independently sourced statement amount could remain usable.
A one-page negative search could never establish absence from the whole filing.

That review proposal has its own issuer and source dependencies. It supplies
no Apple financial constraint, no posting permission and no observed human
decision. The connected financial result and the separate review proposal can
therefore be inspected together without joining their companies' histories.
Section~\ref{sec:real-package-reference} gives the exact exchange locations
and rejection boundaries that preserve these distinctions.

